\documentclass[aspectratio=169,9pt]{beamer}
\usepackage[utf8]{inputenc}
\usepackage[T1]{fontenc}
\usepackage[spanish,mexico]{babel}
\usepackage{amsmath,amssymb,mathtools}
\usepackage{booktabs,array,longtable}
\usepackage{multicol}
\usepackage{microtype}
\usetheme{Madrid}
\usecolortheme{default}
\setbeamertemplate{navigation symbols}{}
\setbeamertemplate{theorems}[numbered]
\setbeamersize{text margin left=7mm,text margin right=7mm}
\newtheorem{definicion}{Definición}[section]
\newtheorem{teorema}[definicion]{Teorema}
\newtheorem{proposicion}[definicion]{Proposición}
\newtheorem{corolario}[definicion]{Corolario}
\theoremstyle{remark}
\newtheorem{observacion}[definicion]{Observación}
\newtheorem{ejemplo}[definicion]{Ejemplo}
\newcommand{\R}{\mathbb{R}}
\newcommand{\rank}{\operatorname{rang}}
\newcommand{\RREF}{\operatorname{RREF}}
\newcommand{\sistema}[1]{%
\[
\left\{
\begin{aligned}
#1
\end{aligned}
\right.
\]
}

\title{Álgebra Lineal}
\subtitle{Unidad 2: Sistemas de ecuaciones lineales}
\author{Carlos Ernesto Martínez}
\institute{Tecnológico Nacional de México}
\date{Agosto--Diciembre 2026}
\begin{document}
\begin{frame}\titlepage\end{frame}
\begin{frame}[allowframebreaks]{Contenido}\tableofcontents[hideallsubsections]\end{frame}
\section{Propósito de la unidad}
\begin{frame}[allowframebreaks,t]{Propósito de la unidad}
En esta unidad estudiaremos los \emph{sistemas de ecuaciones lineales}. El objetivo no es únicamente
aprender algoritmos para obtener valores numéricos. Nos interesa comprender qué significa resolver
simultáneamente varias ecuaciones, cuándo un sistema tiene solución, cuántas soluciones puede tener,
cómo interpretar geométricamente esas posibilidades y por qué los métodos de eliminación funcionan.

La idea central será pasar progresivamente de la escritura usual de un sistema,
\[
\begin{cases}
a_{11}x_1+\cdots+a_{1n}x_n=b_1,\\
\vdots\\
a_{m1}x_1+\cdots+a_{mn}x_n=b_m,
\end{cases}
\]
a una representación matricial compacta
\[
A\mathbf{x}=\mathbf{b}.
\]
La matriz no sustituye al sistema: conserva de manera organizada la información de sus coeficientes y
permite realizar de forma eficiente las mismas operaciones algebraicas que ya conocemos.

\par\medskip{\large\bfseries Ruta conceptual}\par\smallskip
Seguiremos la cadena
\[
\boxed{\text{ecuación lineal}}
\longrightarrow
\boxed{\text{sistema}}
\longrightarrow
\boxed{\text{matriz aumentada}}
\longrightarrow
\boxed{\text{operaciones elementales}}
\longrightarrow
\boxed{\text{Gauss / Gauss--Jordan}}
\]
y después estudiaremos
\[
\boxed{\text{tipos de solución}}
\longrightarrow
\boxed{\text{variables libres}}
\longrightarrow
\boxed{\text{rango y consistencia}}
\longrightarrow
\boxed{\text{métodos alternativos y aplicaciones}}.
\]
\end{frame}

\section{Ecuaciones lineales}
\begin{frame}[allowframebreaks,t]{Ecuaciones lineales}
\begin{definicion}
Una \textbf{ecuación lineal} en las incógnitas $x_1,\ldots,x_n$ es una ecuación de la forma
\[
a_1x_1+a_2x_2+\cdots+a_nx_n=b,
\]
donde $a_1,\ldots,a_n,b\in\R$ son constantes conocidas.
\end{definicion}

La palabra \emph{lineal} significa, en particular, que las incógnitas aparecen únicamente con potencia
uno y no se multiplican entre sí. No aparecen términos como $x^2$, $xy$, $\sqrt{x}$ o $\sin x$.

\begin{ejemplo}
Las expresiones
\[
2x-3y=7,\qquad x_1+4x_2-2x_3=5,\qquad 3x+2y-z=0
\]
son ecuaciones lineales.
\end{ejemplo}

\begin{ejemplo}
Las ecuaciones
\[
x^2+y=3,\qquad xy+z=1,\qquad \sin x+y=0
\]
no son lineales.
\end{ejemplo}

\begin{definicion}
Una \textbf{solución} de una ecuación lineal es un conjunto de valores de las incógnitas que hace verdadera
la igualdad.
\end{definicion}

Por ejemplo, $(x,y)=(2,1)$ es solución de $2x-y=3$, pues
\[
2(2)-1=3.
\]
Sin embargo, una sola ecuación con varias incógnitas suele tener muchas soluciones. El problema se vuelve
más interesante cuando exigimos que los mismos valores satisfagan simultáneamente varias ecuaciones.
\end{frame}

\section{Sistemas de ecuaciones lineales}
\begin{frame}[allowframebreaks,t]{Sistemas de ecuaciones lineales}
\begin{definicion}
Un \textbf{sistema de $m$ ecuaciones lineales con $n$ incógnitas} es un conjunto de ecuaciones
\[
\begin{cases}
a_{11}x_1+a_{12}x_2+\cdots+a_{1n}x_n=b_1,\\
a_{21}x_1+a_{22}x_2+\cdots+a_{2n}x_n=b_2,\\
\hspace{2.5cm}\vdots\\
a_{m1}x_1+a_{m2}x_2+\cdots+a_{mn}x_n=b_m.
\end{cases}
\]
\end{definicion}

Una solución del sistema debe satisfacer \emph{todas} las ecuaciones al mismo tiempo.

\begin{ejemplo}
Consideremos
\[
\begin{cases}
x+y=5,\\
2x-y=1.
\end{cases}
\]
Sumando ambas ecuaciones,
\[
3x=6,
\]
de donde $x=2$. Sustituyendo en la primera ecuación,
\[
2+y=5,
\]
y por tanto $y=3$. La solución es
\[
(x,y)=(2,3).
\]
La verificación es inmediata:
\[
2+3=5,\qquad 2(2)-3=1.
\]
\end{ejemplo}

\par\medskip{\large\bfseries ¿Qué significa resolver un sistema?}\par\smallskip

Resolver un sistema significa describir \emph{todo su conjunto solución}. No basta encontrar una pareja
que funcione si existen otras. Tampoco debemos concluir que ``no se puede resolver'' sólo porque un
procedimiento produce una variable libre: precisamente esa variable libre puede indicar la existencia
de infinitas soluciones.
\end{frame}

\section{Interpretación geométrica}
\begin{frame}[allowframebreaks,t]{Interpretación geométrica}
La geometría permite comprender, antes de desarrollar los algoritmos, por qué aparecen distintos tipos
de solución.

\par\medskip{\large\bfseries Dos ecuaciones con dos incógnitas}\par\smallskip

Una ecuación
\[
ax+by=c
\]
representa, salvo casos degenerados, una recta en $\R^2$. Resolver dos ecuaciones simultáneamente
equivale a buscar los puntos comunes de dos rectas.

Existen tres posibilidades fundamentales.

\par\smallskip{\bfseries Rectas que se cortan}\par
\[
\begin{cases}
x+y=3,\\
x-y=1.
\end{cases}
\]
Las rectas se intersectan en un único punto. El sistema tiene una solución.

\par\smallskip{\bfseries Rectas paralelas distintas}\par
\[
\begin{cases}
x+y=2,\\
2x+2y=5.
\end{cases}
\]
La segunda ecuación tendría el mismo lado izquierdo que el doble de la primera, pero un término
independiente diferente. Las rectas son paralelas y el sistema no tiene solución.

\par\smallskip{\bfseries La misma recta}\par
\[
\begin{cases}
x+y=2,\\
2x+2y=4.
\end{cases}
\]
La segunda ecuación es exactamente el doble de la primera. Ambas describen la misma recta, por lo que
cada punto de ella es solución: existen infinitas soluciones.

\par\medskip{\large\bfseries Tres incógnitas}\par\smallskip

En $\R^3$, una ecuación lineal
\[
ax+by+cz=d
\]
representa normalmente un plano. Resolver un sistema significa buscar la intersección común de los
planos. Pueden cortarse en un punto, en una recta, coincidir parcialmente o no tener ningún punto común.

Esta interpretación anticipa un hecho fundamental.

\begin{proposicion}
Un sistema lineal puede tener:
\begin{enumerate}
\item una única solución;
\item ninguna solución;
\item infinitas soluciones.
\end{enumerate}
No puede tener exactamente dos, tres o cualquier otro número finito mayor que uno de soluciones.
\end{proposicion}

Más adelante podremos justificar formalmente esta afirmación mediante la estructura de los sistemas
homogéneos y las variables libres.
\end{frame}

\section{Representación matricial}
\begin{frame}[allowframebreaks,t]{Representación matricial}
Consideremos
\[
\begin{cases}
2x-y+3z=7,\\
x+4y-z=2,\\
3x+2y+2z=5.
\end{cases}
\]
Separaremos tres objetos:
\[
A=
\begin{pmatrix}
2&-1&3\\
1&4&-1\\
3&2&2
\end{pmatrix},
\qquad
\mathbf{x}=
\begin{pmatrix}x\\y\\z\end{pmatrix},
\qquad
\mathbf{b}=
\begin{pmatrix}7\\2\\5\end{pmatrix}.
\]
Entonces el sistema se escribe
\[
A\mathbf{x}=\mathbf{b}.
\]

\begin{definicion}
$A$ se denomina \textbf{matriz de coeficientes}; $\mathbf{x}$ es el vector de incógnitas y
$\mathbf{b}$ el vector de términos independientes.
\end{definicion}

\begin{definicion}
La \textbf{matriz aumentada} del sistema es
\[
[A\mid\mathbf b]=
\left[
\begin{array}{ccc|c}
2&-1&3&7\\
1&4&-1&2\\
3&2&2&5
\end{array}
\right].
\]
\end{definicion}

La barra vertical es sólo una ayuda visual: separa los coeficientes de las incógnitas del término
independiente.
\end{frame}

\section{Sistemas homogéneos y no homogéneos}
\begin{frame}[allowframebreaks,t]{Sistemas homogéneos y no homogéneos}
\begin{definicion}
Un sistema es \textbf{homogéneo} si todos sus términos independientes son cero:
\[
A\mathbf{x}=\mathbf{0}.
\]
Es \textbf{no homogéneo} si
\[
A\mathbf{x}=\mathbf{b},\qquad \mathbf b\neq\mathbf0.
\]
\end{definicion}

Todo sistema homogéneo posee al menos una solución:
\[
\mathbf{x}=\mathbf0.
\]
Ésta se denomina \textbf{solución trivial}.

\begin{ejemplo}
\[
\begin{cases}
x+2y-z=0,\\
2x+4y-2z=0.
\end{cases}
\]
La segunda ecuación es el doble de la primera. Tomemos
\[
y=s,\qquad z=t.
\]
Entonces
\[
x=-2s+t,
\]
y todas las soluciones son
\[
\begin{pmatrix}x\\y\\z\end{pmatrix}
=
\begin{pmatrix}-2s+t\\s\\t\end{pmatrix}
=
s\begin{pmatrix}-2\\1\\0\end{pmatrix}
+t\begin{pmatrix}1\\0\\1\end{pmatrix}.
\]
Además de la solución trivial $s=t=0$, existen infinitas soluciones no triviales.
\end{ejemplo}

Esta forma de escribir las soluciones será especialmente importante en unidades posteriores: ya
aparece, de manera natural, la idea de combinación lineal.
\end{frame}

\section{Sistemas compatibles e incompatibles}
\begin{frame}[allowframebreaks,t]{Sistemas compatibles e incompatibles}
\begin{definicion}
Un sistema es \textbf{compatible} si tiene al menos una solución. Es \textbf{incompatible} si no tiene
solución.
\end{definicion}

Un sistema compatible puede ser:
\begin{itemize}
\item \textbf{compatible determinado}: tiene una única solución;
\item \textbf{compatible indeterminado}: tiene infinitas soluciones.
\end{itemize}

Por tanto,
\[
\boxed{
\begin{array}{c}
\text{sistema lineal}\\[1mm]
\Downarrow\\
\text{compatible}\quad\text{o}\quad\text{incompatible}\\
\text{una o infinitas soluciones}\qquad\text{ninguna solución}
\end{array}}
\]
\end{frame}

\section{De la eliminación algebraica a las operaciones por renglones}
\begin{frame}[allowframebreaks,t]{De la eliminación algebraica a las operaciones por renglones}
Antes de introducir un algoritmo matricial, recordemos la eliminación usual:
\[
\begin{cases}
x+2y=5,\\
3x-y=4.
\end{cases}
\]
Para eliminar $x$, reemplazamos la segunda ecuación por
\[
E_2-3E_1.
\]
Obtenemos
\[
-7y=-11.
\]
Lo importante es observar que no hemos cambiado el conjunto solución: sólo reemplazamos una ecuación
por otra consecuencia algebraicamente equivalente.

La matriz aumentada es
\[
\left[
\begin{array}{cc|c}
1&2&5\\
3&-1&4
\end{array}
\right].
\]
La misma operación se escribe
\[
R_2\leftarrow R_2-3R_1,
\]
y produce
\[
\left[
\begin{array}{cc|c}
1&2&5\\
0&-7&-11
\end{array}
\right].
\]

\par\medskip{\large\bfseries Operaciones elementales}\par\smallskip

\begin{definicion}
Las \textbf{operaciones elementales por renglones} son:
\begin{enumerate}
\item intercambiar dos renglones:
\[
R_i\leftrightarrow R_j;
\]
\item multiplicar un renglón por un escalar no nulo:
\[
R_i\leftarrow cR_i,\qquad c\neq0;
\]
\item sumar a un renglón un múltiplo de otro:
\[
R_i\leftarrow R_i+cR_j.
\]
\end{enumerate}
\end{definicion}

\begin{teorema}
Las operaciones elementales por renglones preservan el conjunto solución del sistema.
\end{teorema}

\begin{proof}
Intercambiar ecuaciones no modifica la exigencia de que todas se satisfagan. Multiplicar una ecuación
por un número no nulo produce una ecuación equivalente. Finalmente, sustituir una ecuación por ella
misma más un múltiplo de otra es reversible:
\[
R_i\leftarrow R_i+cR_j
\quad\Longleftrightarrow\quad
R_i\leftarrow R_i-cR_j.
\]
Por tanto, cada operación transforma el sistema en otro equivalente.
\end{proof}
\end{frame}

\section{Forma escalonada}
\begin{frame}[allowframebreaks,t]{Forma escalonada}
\begin{definicion}
Una matriz está en \textbf{forma escalonada por renglones} si:
\begin{enumerate}
\item todos los renglones nulos están debajo de los no nulos;
\item el primer elemento no nulo de cada renglón aparece a la derecha del primer elemento no nulo
del renglón anterior;
\item debajo de cada pivote hay ceros.
\end{enumerate}
\end{definicion}

Por ejemplo,
\[
\begin{pmatrix}
1&2&-1&3\\
0&3&4&2\\
0&0&5&-1\\
0&0&0&0
\end{pmatrix}
\]
está en forma escalonada.

\begin{definicion}
La primera entrada no nula de un renglón no nulo se denomina \textbf{pivote}. La columna que contiene
un pivote se denomina \textbf{columna pivote}.
\end{definicion}
\end{frame}

\section{Eliminación de Gauss}
\begin{frame}[allowframebreaks,t]{Eliminación de Gauss}
El método de Gauss consiste en aplicar operaciones elementales hasta llevar la matriz aumentada a forma
escalonada. Después se utiliza sustitución hacia atrás.

\par\medskip{\large\bfseries Ejemplo con solución única}\par\smallskip

Resolvamos
\[
\begin{cases}
x+y+z=6,\\
2x-y+z=3,\\
x+2y-z=2.
\end{cases}
\]
La matriz aumentada es
\[
\left[
\begin{array}{ccc|c}
1&1&1&6\\
2&-1&1&3\\
1&2&-1&2
\end{array}
\right].
\]
Aplicamos
\[
R_2\leftarrow R_2-2R_1,\qquad
R_3\leftarrow R_3-R_1,
\]
para obtener
\[
\left[
\begin{array}{ccc|c}
1&1&1&6\\
0&-3&-1&-9\\
0&1&-2&-4
\end{array}
\right].
\]
Conviene intercambiar los últimos dos renglones:
\[
R_2\leftrightarrow R_3,
\]
de modo que
\[
\left[
\begin{array}{ccc|c}
1&1&1&6\\
0&1&-2&-4\\
0&-3&-1&-9
\end{array}
\right].
\]
Ahora
\[
R_3\leftarrow R_3+3R_2,
\]
y obtenemos
\[
\left[
\begin{array}{ccc|c}
1&1&1&6\\
0&1&-2&-4\\
0&0&-7&-21
\end{array}
\right].
\]
La última ecuación da $z=3$. Luego
\[
y-2(3)=-4\quad\Longrightarrow\quad y=2,
\]
y finalmente
\[
x+2+3=6\quad\Longrightarrow\quad x=1.
\]
Así,
\[
\boxed{(x,y,z)=(1,2,3)}.
\]

\par\medskip{\large\bfseries Ejemplo incompatible}\par\smallskip

\[
\begin{cases}
x+y=2,\\
2x+2y=5.
\end{cases}
\]
Tenemos
\[
\left[
\begin{array}{cc|c}
1&1&2\\
2&2&5
\end{array}
\right]
\xrightarrow{R_2-2R_1}
\left[
\begin{array}{cc|c}
1&1&2\\
0&0&1
\end{array}
\right].
\]
El último renglón representa
\[
0x+0y=1,
\]
es decir,
\[
0=1,
\]
una contradicción. Por tanto, el sistema es incompatible.

\par\medskip{\large\bfseries Ejemplo con infinitas soluciones}\par\smallskip

\[
\begin{cases}
x+y+z=2,\\
2x+2y+2z=4.
\end{cases}
\]
La reducción produce
\[
\left[
\begin{array}{ccc|c}
1&1&1&2\\
0&0&0&0
\end{array}
\right].
\]
Sólo existe una ecuación independiente para tres incógnitas. Tomemos
\[
y=s,\qquad z=t.
\]
Entonces
\[
x=2-s-t.
\]
Por tanto,
\[
\boxed{(x,y,z)=(2-s-t,s,t),\qquad s,t\in\R.}
\]
\end{frame}

\section{Variables básicas y variables libres}
\begin{frame}[allowframebreaks,t]{Variables básicas y variables libres}
En una matriz escalonada, las variables correspondientes a columnas pivote se denominan
\textbf{variables básicas}. Las restantes son \textbf{variables libres}.

En
\[
\left[
\begin{array}{cccc|c}
1&2&0&-1&3\\
0&0&1&4&2
\end{array}
\right],
\]
las columnas pivote son la primera y la tercera. Por tanto,
\[
x_1,\ x_3
\]
son básicas, mientras que
\[
x_2,\ x_4
\]
son libres.

Tomando
\[
x_2=s,\qquad x_4=t,
\]
obtenemos
\[
x_1=3-2s+t,\qquad x_3=2-4t.
\]
La solución completa es
\[
\mathbf{x}=
\begin{pmatrix}3\\0\\2\\0\end{pmatrix}
+s\begin{pmatrix}-2\\1\\0\\0\end{pmatrix}
+t\begin{pmatrix}1\\0\\-4\\1\end{pmatrix}.
\]
\end{frame}

\section{Forma escalonada reducida}
\begin{frame}[allowframebreaks,t]{Forma escalonada reducida}
\begin{definicion}
Una matriz está en \textbf{forma escalonada reducida por renglones} si está en forma escalonada y,
además:
\begin{enumerate}
\item cada pivote es igual a $1$;
\item cada pivote es el único elemento no nulo de su columna.
\end{enumerate}
\end{definicion}

Por ejemplo,
\[
\begin{pmatrix}
1&0&2&0&3\\
0&1&-1&0&4\\
0&0&0&1&-2
\end{pmatrix}
\]
está en forma escalonada reducida.

\begin{teorema}
Toda matriz es equivalente por renglones a una única matriz escalonada reducida.
\end{teorema}

Esta unicidad es importante: podemos realizar operaciones distintas durante el proceso, pero si
continuamos hasta la forma reducida, el resultado final será el mismo.
\end{frame}

\section{Método de Gauss--Jordan}
\begin{frame}[t]{Gauss--Jordan: idea del algoritmo}
Gauss--Jordan continúa la eliminación hasta obtener RREF. En cada etapa preguntamos:
\[
\boxed{\text{¿Cuál es el pivote actual y qué entradas de su columna debo anular?}}
\]
\begin{enumerate}
\item Escribir la matriz aumentada.
\item Elegir la columna pivote; intercambiar renglones si es necesario.
\item Normalizar el pivote a \(1\).
\item Hacer ceros debajo y encima del pivote.
\item Repetir en la siguiente columna.
\item Interpretar la RREF y comprobar.
\end{enumerate}
\end{frame}
\subsection{Ejemplo $2\times2$}
\begin{frame}[t]{Ejemplo $2\times2$: planteamiento}
\textbf{Objetivo.} Mecánica elemental.
\[
\begin{aligned}
x_{1} - x_{2} - 3 x_{4} &= -7\\
- 2 x_{1} + 2 x_{2} + 3 x_{3} - 3 x_{4} &= 5\\
- x_{1} - 3 x_{3} - 2 x_{4} - 2 x_{5} &= -20\\
- 2 x_{1} - 3 x_{2} - 3 x_{3} + x_{4} - 2 x_{5} &= -21\\
- 2 x_{1} - 3 x_{2} - 3 x_{3} + x_{4} - 2 x_{5} &= -20
\end{aligned}
\]
Matriz aumentada:
\[
\left[\begin{array}{ccccc|c}
1 & -1 & 0 & -3 & 0 & -7\\
-2 & 2 & 3 & -3 & 0 & 5\\
-1 & 0 & -3 & -2 & -2 & -20\\
-2 & -3 & -3 & 1 & -2 & -21\\
-2 & -3 & -3 & 1 & -2 & -20
\end{array}\right]
\]
\end{frame}
\begin{frame}[t]{Ejemplo $2\times2$: paso 1}
Usamos el pivote de la columna 1 para anular la entrada del renglón 2.
\[
R_{2}\leftarrow R_{2}+2R_{1}
\]
\[
\left[\begin{array}{ccccc|c}
1 & -1 & 0 & -3 & 0 & -7\\
0 & 0 & 3 & -9 & 0 & -9\\
-1 & 0 & -3 & -2 & -2 & -20\\
-2 & -3 & -3 & 1 & -2 & -21\\
-2 & -3 & -3 & 1 & -2 & -20
\end{array}\right]
\]
\end{frame}
\begin{frame}[t]{Ejemplo $2\times2$: paso 2}
Usamos el pivote de la columna 1 para anular la entrada del renglón 3.
\[
R_{3}\leftarrow R_{3}+R_{1}
\]
\[
\left[\begin{array}{ccccc|c}
1 & -1 & 0 & -3 & 0 & -7\\
0 & 0 & 3 & -9 & 0 & -9\\
0 & -1 & -3 & -5 & -2 & -27\\
-2 & -3 & -3 & 1 & -2 & -21\\
-2 & -3 & -3 & 1 & -2 & -20
\end{array}\right]
\]
\end{frame}
\begin{frame}[t]{Ejemplo $2\times2$: paso 3}
Usamos el pivote de la columna 1 para anular la entrada del renglón 4.
\[
R_{4}\leftarrow R_{4}+2R_{1}
\]
\[
\left[\begin{array}{ccccc|c}
1 & -1 & 0 & -3 & 0 & -7\\
0 & 0 & 3 & -9 & 0 & -9\\
0 & -1 & -3 & -5 & -2 & -27\\
0 & -5 & -3 & -5 & -2 & -35\\
-2 & -3 & -3 & 1 & -2 & -20
\end{array}\right]
\]
\end{frame}
\begin{frame}[t]{Ejemplo $2\times2$: paso 4}
Usamos el pivote de la columna 1 para anular la entrada del renglón 5.
\[
R_{5}\leftarrow R_{5}+2R_{1}
\]
\[
\left[\begin{array}{ccccc|c}
1 & -1 & 0 & -3 & 0 & -7\\
0 & 0 & 3 & -9 & 0 & -9\\
0 & -1 & -3 & -5 & -2 & -27\\
0 & -5 & -3 & -5 & -2 & -35\\
0 & -5 & -3 & -5 & -2 & -34
\end{array}\right]
\]
\end{frame}
\begin{frame}[t]{Ejemplo $2\times2$: paso 5}
Intercambiamos renglones para colocar un elemento no nulo en la posición pivote de la columna 2.
\[
R_{2}\leftrightarrow R_{3}
\]
\[
\left[\begin{array}{ccccc|c}
1 & -1 & 0 & -3 & 0 & -7\\
0 & -1 & -3 & -5 & -2 & -27\\
0 & 0 & 3 & -9 & 0 & -9\\
0 & -5 & -3 & -5 & -2 & -35\\
0 & -5 & -3 & -5 & -2 & -34
\end{array}\right]
\]
\end{frame}
\begin{frame}[t]{Ejemplo $2\times2$: paso 6}
Normalizamos el pivote de la columna 2 para convertirlo en 1.
\[
R_{2}\leftarrow \frac{1}{-1}R_{2}
\]
\[
\left[\begin{array}{ccccc|c}
1 & -1 & 0 & -3 & 0 & -7\\
0 & 1 & 3 & 5 & 2 & 27\\
0 & 0 & 3 & -9 & 0 & -9\\
0 & -5 & -3 & -5 & -2 & -35\\
0 & -5 & -3 & -5 & -2 & -34
\end{array}\right]
\]
\end{frame}
\begin{frame}[t]{Ejemplo $2\times2$: paso 7}
Usamos el pivote de la columna 2 para anular la entrada del renglón 1.
\[
R_{1}\leftarrow R_{1}+R_{2}
\]
\[
\left[\begin{array}{ccccc|c}
1 & 0 & 3 & 2 & 2 & 20\\
0 & 1 & 3 & 5 & 2 & 27\\
0 & 0 & 3 & -9 & 0 & -9\\
0 & -5 & -3 & -5 & -2 & -35\\
0 & -5 & -3 & -5 & -2 & -34
\end{array}\right]
\]
\end{frame}
\begin{frame}[t]{Ejemplo $2\times2$: paso 8}
Usamos el pivote de la columna 2 para anular la entrada del renglón 4.
\[
R_{4}\leftarrow R_{4}+5R_{2}
\]
\[
\left[\begin{array}{ccccc|c}
1 & 0 & 3 & 2 & 2 & 20\\
0 & 1 & 3 & 5 & 2 & 27\\
0 & 0 & 3 & -9 & 0 & -9\\
0 & 0 & 12 & 20 & 8 & 100\\
0 & -5 & -3 & -5 & -2 & -34
\end{array}\right]
\]
\end{frame}
\begin{frame}[t]{Ejemplo $2\times2$: paso 9}
Usamos el pivote de la columna 2 para anular la entrada del renglón 5.
\[
R_{5}\leftarrow R_{5}+5R_{2}
\]
\[
\left[\begin{array}{ccccc|c}
1 & 0 & 3 & 2 & 2 & 20\\
0 & 1 & 3 & 5 & 2 & 27\\
0 & 0 & 3 & -9 & 0 & -9\\
0 & 0 & 12 & 20 & 8 & 100\\
0 & 0 & 12 & 20 & 8 & 101
\end{array}\right]
\]
\end{frame}
\begin{frame}[t]{Ejemplo $2\times2$: paso 10}
Normalizamos el pivote de la columna 3 para convertirlo en 1.
\[
R_{3}\leftarrow \frac{1}{3}R_{3}
\]
\[
\left[\begin{array}{ccccc|c}
1 & 0 & 3 & 2 & 2 & 20\\
0 & 1 & 3 & 5 & 2 & 27\\
0 & 0 & 1 & -3 & 0 & -3\\
0 & 0 & 12 & 20 & 8 & 100\\
0 & 0 & 12 & 20 & 8 & 101
\end{array}\right]
\]
\end{frame}
\begin{frame}[t]{Ejemplo $2\times2$: paso 11}
Usamos el pivote de la columna 3 para anular la entrada del renglón 1.
\[
R_{1}\leftarrow R_{1}-3R_{3}
\]
\[
\left[\begin{array}{ccccc|c}
1 & 0 & 0 & 11 & 2 & 29\\
0 & 1 & 3 & 5 & 2 & 27\\
0 & 0 & 1 & -3 & 0 & -3\\
0 & 0 & 12 & 20 & 8 & 100\\
0 & 0 & 12 & 20 & 8 & 101
\end{array}\right]
\]
\end{frame}
\begin{frame}[t]{Ejemplo $2\times2$: paso 12}
Usamos el pivote de la columna 3 para anular la entrada del renglón 2.
\[
R_{2}\leftarrow R_{2}-3R_{3}
\]
\[
\left[\begin{array}{ccccc|c}
1 & 0 & 0 & 11 & 2 & 29\\
0 & 1 & 0 & 14 & 2 & 36\\
0 & 0 & 1 & -3 & 0 & -3\\
0 & 0 & 12 & 20 & 8 & 100\\
0 & 0 & 12 & 20 & 8 & 101
\end{array}\right]
\]
\end{frame}
\begin{frame}[t]{Ejemplo $2\times2$: paso 13}
Usamos el pivote de la columna 3 para anular la entrada del renglón 4.
\[
R_{4}\leftarrow R_{4}-12R_{3}
\]
\[
\left[\begin{array}{ccccc|c}
1 & 0 & 0 & 11 & 2 & 29\\
0 & 1 & 0 & 14 & 2 & 36\\
0 & 0 & 1 & -3 & 0 & -3\\
0 & 0 & 0 & 56 & 8 & 136\\
0 & 0 & 12 & 20 & 8 & 101
\end{array}\right]
\]
\end{frame}
\begin{frame}[t]{Ejemplo $2\times2$: paso 14}
Usamos el pivote de la columna 3 para anular la entrada del renglón 5.
\[
R_{5}\leftarrow R_{5}-12R_{3}
\]
\[
\left[\begin{array}{ccccc|c}
1 & 0 & 0 & 11 & 2 & 29\\
0 & 1 & 0 & 14 & 2 & 36\\
0 & 0 & 1 & -3 & 0 & -3\\
0 & 0 & 0 & 56 & 8 & 136\\
0 & 0 & 0 & 56 & 8 & 137
\end{array}\right]
\]
\end{frame}
\begin{frame}[t]{Ejemplo $2\times2$: paso 15}
Normalizamos el pivote de la columna 4 para convertirlo en 1.
\[
R_{4}\leftarrow \frac{1}{56}R_{4}
\]
\[
\left[\begin{array}{ccccc|c}
1 & 0 & 0 & 11 & 2 & 29\\
0 & 1 & 0 & 14 & 2 & 36\\
0 & 0 & 1 & -3 & 0 & -3\\
0 & 0 & 0 & 1 & \frac{1}{7} & \frac{17}{7}\\
0 & 0 & 0 & 56 & 8 & 137
\end{array}\right]
\]
\end{frame}
\begin{frame}[t]{Ejemplo $2\times2$: paso 16}
Usamos el pivote de la columna 4 para anular la entrada del renglón 1.
\[
R_{1}\leftarrow R_{1}-11R_{4}
\]
\[
\left[\begin{array}{ccccc|c}
1 & 0 & 0 & 0 & \frac{3}{7} & \frac{16}{7}\\
0 & 1 & 0 & 14 & 2 & 36\\
0 & 0 & 1 & -3 & 0 & -3\\
0 & 0 & 0 & 1 & \frac{1}{7} & \frac{17}{7}\\
0 & 0 & 0 & 56 & 8 & 137
\end{array}\right]
\]
\end{frame}
\begin{frame}[t]{Ejemplo $2\times2$: paso 17}
Usamos el pivote de la columna 4 para anular la entrada del renglón 2.
\[
R_{2}\leftarrow R_{2}-14R_{4}
\]
\[
\left[\begin{array}{ccccc|c}
1 & 0 & 0 & 0 & \frac{3}{7} & \frac{16}{7}\\
0 & 1 & 0 & 0 & 0 & 2\\
0 & 0 & 1 & -3 & 0 & -3\\
0 & 0 & 0 & 1 & \frac{1}{7} & \frac{17}{7}\\
0 & 0 & 0 & 56 & 8 & 137
\end{array}\right]
\]
\end{frame}
\begin{frame}[t]{Ejemplo $2\times2$: paso 18}
Usamos el pivote de la columna 4 para anular la entrada del renglón 3.
\[
R_{3}\leftarrow R_{3}+3R_{4}
\]
\[
\left[\begin{array}{ccccc|c}
1 & 0 & 0 & 0 & \frac{3}{7} & \frac{16}{7}\\
0 & 1 & 0 & 0 & 0 & 2\\
0 & 0 & 1 & 0 & \frac{3}{7} & \frac{30}{7}\\
0 & 0 & 0 & 1 & \frac{1}{7} & \frac{17}{7}\\
0 & 0 & 0 & 56 & 8 & 137
\end{array}\right]
\]
\end{frame}
\begin{frame}[t]{Ejemplo $2\times2$: paso 19}
Usamos el pivote de la columna 4 para anular la entrada del renglón 5.
\[
R_{5}\leftarrow R_{5}-56R_{4}
\]
\[
\left[\begin{array}{ccccc|c}
1 & 0 & 0 & 0 & \frac{3}{7} & \frac{16}{7}\\
0 & 1 & 0 & 0 & 0 & 2\\
0 & 0 & 1 & 0 & \frac{3}{7} & \frac{30}{7}\\
0 & 0 & 0 & 1 & \frac{1}{7} & \frac{17}{7}\\
0 & 0 & 0 & 0 & 0 & 1
\end{array}\right]
\]
\end{frame}
\begin{frame}[t]{Ejemplo $2\times2$: conclusión}
La matriz está en forma escalonada reducida. La solución se lee directamente:
\[
\boxed{x_{1}=\frac{16}{7},\quad x_{2}=2,\quad x_{3}=\frac{30}{7},\quad x_{4}=\frac{17}{7},\quad x_{5}=1}
\]
\begin{block}{Comprobación}
Sustituimos los valores en el sistema original. Si se recuperan todos los términos independientes, la solución es correcta.
\end{block}
\end{frame}
\subsection{Ejemplo $3\times3$}
\begin{frame}[t]{Ejemplo $3\times3$: planteamiento}
\textbf{Objetivo.} Algoritmo completo.
\[
\begin{aligned}
x_{1} + x_{2} + x_{3} &= 6\\
2 x_{1} - x_{2} + x_{3} &= 3\\
x_{1} + 2 x_{2} - x_{3} &= 2
\end{aligned}
\]
Matriz aumentada:
\[
\left[\begin{array}{ccc|c}
1 & 1 & 1 & 6\\
2 & -1 & 1 & 3\\
1 & 2 & -1 & 2
\end{array}\right]
\]
\end{frame}
\begin{frame}[t]{Ejemplo $3\times3$: paso 1}
Usamos el pivote de la columna 1 para anular la entrada del renglón 2.
\[
R_{2}\leftarrow R_{2}-2R_{1}
\]
\[
\left[\begin{array}{ccc|c}
1 & 1 & 1 & 6\\
0 & -3 & -1 & -9\\
1 & 2 & -1 & 2
\end{array}\right]
\]
\end{frame}
\begin{frame}[t]{Ejemplo $3\times3$: paso 2}
Usamos el pivote de la columna 1 para anular la entrada del renglón 3.
\[
R_{3}\leftarrow R_{3}-R_{1}
\]
\[
\left[\begin{array}{ccc|c}
1 & 1 & 1 & 6\\
0 & -3 & -1 & -9\\
0 & 1 & -2 & -4
\end{array}\right]
\]
\end{frame}
\begin{frame}[t]{Ejemplo $3\times3$: paso 3}
Normalizamos el pivote de la columna 2 para convertirlo en 1.
\[
R_{2}\leftarrow \frac{1}{-3}R_{2}
\]
\[
\left[\begin{array}{ccc|c}
1 & 1 & 1 & 6\\
0 & 1 & \frac{1}{3} & 3\\
0 & 1 & -2 & -4
\end{array}\right]
\]
\end{frame}
\begin{frame}[t]{Ejemplo $3\times3$: paso 4}
Usamos el pivote de la columna 2 para anular la entrada del renglón 1.
\[
R_{1}\leftarrow R_{1}-R_{2}
\]
\[
\left[\begin{array}{ccc|c}
1 & 0 & \frac{2}{3} & 3\\
0 & 1 & \frac{1}{3} & 3\\
0 & 1 & -2 & -4
\end{array}\right]
\]
\end{frame}
\begin{frame}[t]{Ejemplo $3\times3$: paso 5}
Usamos el pivote de la columna 2 para anular la entrada del renglón 3.
\[
R_{3}\leftarrow R_{3}-R_{2}
\]
\[
\left[\begin{array}{ccc|c}
1 & 0 & \frac{2}{3} & 3\\
0 & 1 & \frac{1}{3} & 3\\
0 & 0 & - \frac{7}{3} & -7
\end{array}\right]
\]
\end{frame}
\begin{frame}[t]{Ejemplo $3\times3$: paso 6}
Normalizamos el pivote de la columna 3 para convertirlo en 1.
\[
R_{3}\leftarrow \frac{1}{- \frac{7}{3}}R_{3}
\]
\[
\left[\begin{array}{ccc|c}
1 & 0 & \frac{2}{3} & 3\\
0 & 1 & \frac{1}{3} & 3\\
0 & 0 & 1 & 3
\end{array}\right]
\]
\end{frame}
\begin{frame}[t]{Ejemplo $3\times3$: paso 7}
Usamos el pivote de la columna 3 para anular la entrada del renglón 1.
\[
R_{1}\leftarrow R_{1}-\frac{2}{3}R_{3}
\]
\[
\left[\begin{array}{ccc|c}
1 & 0 & 0 & 1\\
0 & 1 & \frac{1}{3} & 3\\
0 & 0 & 1 & 3
\end{array}\right]
\]
\end{frame}
\begin{frame}[t]{Ejemplo $3\times3$: paso 8}
Usamos el pivote de la columna 3 para anular la entrada del renglón 2.
\[
R_{2}\leftarrow R_{2}-\frac{1}{3}R_{3}
\]
\[
\left[\begin{array}{ccc|c}
1 & 0 & 0 & 1\\
0 & 1 & 0 & 2\\
0 & 0 & 1 & 3
\end{array}\right]
\]
\end{frame}
\begin{frame}[t]{Ejemplo $3\times3$: conclusión}
La matriz está en forma escalonada reducida. La solución se lee directamente:
\[
\boxed{x_{1}=1,\quad x_{2}=2,\quad x_{3}=3}
\]
\begin{block}{Comprobación}
Sustituimos los valores en el sistema original. Si se recuperan todos los términos independientes, la solución es correcta.
\end{block}
\end{frame}
\subsection{Ejemplo con intercambio}
\begin{frame}[t]{Ejemplo con intercambio: planteamiento}
\textbf{Objetivo.} Resolver un pivote inicial nulo.
\[
\begin{aligned}
x_{2} + x_{3} &= 2\\
x_{1} - x_{2} + 2 x_{3} &= 9\\
2 x_{1} + x_{2} - x_{3} &= 0
\end{aligned}
\]
Matriz aumentada:
\[
\left[\begin{array}{ccc|c}
0 & 1 & 1 & 2\\
1 & -1 & 2 & 9\\
2 & 1 & -1 & 0
\end{array}\right]
\]
\end{frame}
\begin{frame}[t]{Ejemplo con intercambio: paso 1}
Intercambiamos renglones para colocar un elemento no nulo en la posición pivote de la columna 1.
\[
R_{1}\leftrightarrow R_{2}
\]
\[
\left[\begin{array}{ccc|c}
1 & -1 & 2 & 9\\
0 & 1 & 1 & 2\\
2 & 1 & -1 & 0
\end{array}\right]
\]
\end{frame}
\begin{frame}[t]{Ejemplo con intercambio: paso 2}
Usamos el pivote de la columna 1 para anular la entrada del renglón 3.
\[
R_{3}\leftarrow R_{3}-2R_{1}
\]
\[
\left[\begin{array}{ccc|c}
1 & -1 & 2 & 9\\
0 & 1 & 1 & 2\\
0 & 3 & -5 & -18
\end{array}\right]
\]
\end{frame}
\begin{frame}[t]{Ejemplo con intercambio: paso 3}
Usamos el pivote de la columna 2 para anular la entrada del renglón 1.
\[
R_{1}\leftarrow R_{1}+R_{2}
\]
\[
\left[\begin{array}{ccc|c}
1 & 0 & 3 & 11\\
0 & 1 & 1 & 2\\
0 & 3 & -5 & -18
\end{array}\right]
\]
\end{frame}
\begin{frame}[t]{Ejemplo con intercambio: paso 4}
Usamos el pivote de la columna 2 para anular la entrada del renglón 3.
\[
R_{3}\leftarrow R_{3}-3R_{2}
\]
\[
\left[\begin{array}{ccc|c}
1 & 0 & 3 & 11\\
0 & 1 & 1 & 2\\
0 & 0 & -8 & -24
\end{array}\right]
\]
\end{frame}
\begin{frame}[t]{Ejemplo con intercambio: paso 5}
Normalizamos el pivote de la columna 3 para convertirlo en 1.
\[
R_{3}\leftarrow \frac{1}{-8}R_{3}
\]
\[
\left[\begin{array}{ccc|c}
1 & 0 & 3 & 11\\
0 & 1 & 1 & 2\\
0 & 0 & 1 & 3
\end{array}\right]
\]
\end{frame}
\begin{frame}[t]{Ejemplo con intercambio: paso 6}
Usamos el pivote de la columna 3 para anular la entrada del renglón 1.
\[
R_{1}\leftarrow R_{1}-3R_{3}
\]
\[
\left[\begin{array}{ccc|c}
1 & 0 & 0 & 2\\
0 & 1 & 1 & 2\\
0 & 0 & 1 & 3
\end{array}\right]
\]
\end{frame}
\begin{frame}[t]{Ejemplo con intercambio: paso 7}
Usamos el pivote de la columna 3 para anular la entrada del renglón 2.
\[
R_{2}\leftarrow R_{2}-R_{3}
\]
\[
\left[\begin{array}{ccc|c}
1 & 0 & 0 & 2\\
0 & 1 & 0 & -1\\
0 & 0 & 1 & 3
\end{array}\right]
\]
\end{frame}
\begin{frame}[t]{Ejemplo con intercambio: conclusión}
La matriz está en forma escalonada reducida. La solución se lee directamente:
\[
\boxed{x_{1}=2,\quad x_{2}=-1,\quad x_{3}=3}
\]
\begin{block}{Comprobación}
Sustituimos los valores en el sistema original. Si se recuperan todos los términos independientes, la solución es correcta.
\end{block}
\end{frame}
\subsection{Ejemplo $4\times4$}
\begin{frame}[t]{Ejemplo $4\times4$: planteamiento}
\textbf{Objetivo.} Organización sistemática.
\[
\begin{aligned}
x_{1} + x_{2} - x_{4} &= -3\\
2 x_{1} - x_{2} + x_{3} + x_{4} &= 8\\
x_{1} + 2 x_{2} - x_{3} + 2 x_{4} &= 3\\
3 x_{1} + 2 x_{3} - x_{4} &= 4
\end{aligned}
\]
Matriz aumentada:
\[
\left[\begin{array}{cccc|c}
1 & 1 & 0 & -1 & -3\\
2 & -1 & 1 & 1 & 8\\
1 & 2 & -1 & 2 & 3\\
3 & 0 & 2 & -1 & 4
\end{array}\right]
\]
\end{frame}
\begin{frame}[t]{Ejemplo $4\times4$: paso 1}
Usamos el pivote de la columna 1 para anular la entrada del renglón 2.
\[
R_{2}\leftarrow R_{2}-2R_{1}
\]
\[
\left[\begin{array}{cccc|c}
1 & 1 & 0 & -1 & -3\\
0 & -3 & 1 & 3 & 14\\
1 & 2 & -1 & 2 & 3\\
3 & 0 & 2 & -1 & 4
\end{array}\right]
\]
\end{frame}
\begin{frame}[t]{Ejemplo $4\times4$: paso 2}
Usamos el pivote de la columna 1 para anular la entrada del renglón 3.
\[
R_{3}\leftarrow R_{3}-R_{1}
\]
\[
\left[\begin{array}{cccc|c}
1 & 1 & 0 & -1 & -3\\
0 & -3 & 1 & 3 & 14\\
0 & 1 & -1 & 3 & 6\\
3 & 0 & 2 & -1 & 4
\end{array}\right]
\]
\end{frame}
\begin{frame}[t]{Ejemplo $4\times4$: paso 3}
Usamos el pivote de la columna 1 para anular la entrada del renglón 4.
\[
R_{4}\leftarrow R_{4}-3R_{1}
\]
\[
\left[\begin{array}{cccc|c}
1 & 1 & 0 & -1 & -3\\
0 & -3 & 1 & 3 & 14\\
0 & 1 & -1 & 3 & 6\\
0 & -3 & 2 & 2 & 13
\end{array}\right]
\]
\end{frame}
\begin{frame}[t]{Ejemplo $4\times4$: paso 4}
Normalizamos el pivote de la columna 2 para convertirlo en 1.
\[
R_{2}\leftarrow \frac{1}{-3}R_{2}
\]
\[
\left[\begin{array}{cccc|c}
1 & 1 & 0 & -1 & -3\\
0 & 1 & - \frac{1}{3} & -1 & - \frac{14}{3}\\
0 & 1 & -1 & 3 & 6\\
0 & -3 & 2 & 2 & 13
\end{array}\right]
\]
\end{frame}
\begin{frame}[t]{Ejemplo $4\times4$: paso 5}
Usamos el pivote de la columna 2 para anular la entrada del renglón 1.
\[
R_{1}\leftarrow R_{1}-R_{2}
\]
\[
\left[\begin{array}{cccc|c}
1 & 0 & \frac{1}{3} & 0 & \frac{5}{3}\\
0 & 1 & - \frac{1}{3} & -1 & - \frac{14}{3}\\
0 & 1 & -1 & 3 & 6\\
0 & -3 & 2 & 2 & 13
\end{array}\right]
\]
\end{frame}
\begin{frame}[t]{Ejemplo $4\times4$: paso 6}
Usamos el pivote de la columna 2 para anular la entrada del renglón 3.
\[
R_{3}\leftarrow R_{3}-R_{2}
\]
\[
\left[\begin{array}{cccc|c}
1 & 0 & \frac{1}{3} & 0 & \frac{5}{3}\\
0 & 1 & - \frac{1}{3} & -1 & - \frac{14}{3}\\
0 & 0 & - \frac{2}{3} & 4 & \frac{32}{3}\\
0 & -3 & 2 & 2 & 13
\end{array}\right]
\]
\end{frame}
\begin{frame}[t]{Ejemplo $4\times4$: paso 7}
Usamos el pivote de la columna 2 para anular la entrada del renglón 4.
\[
R_{4}\leftarrow R_{4}+3R_{2}
\]
\[
\left[\begin{array}{cccc|c}
1 & 0 & \frac{1}{3} & 0 & \frac{5}{3}\\
0 & 1 & - \frac{1}{3} & -1 & - \frac{14}{3}\\
0 & 0 & - \frac{2}{3} & 4 & \frac{32}{3}\\
0 & 0 & 1 & -1 & -1
\end{array}\right]
\]
\end{frame}
\begin{frame}[t]{Ejemplo $4\times4$: paso 8}
Normalizamos el pivote de la columna 3 para convertirlo en 1.
\[
R_{3}\leftarrow \frac{1}{- \frac{2}{3}}R_{3}
\]
\[
\left[\begin{array}{cccc|c}
1 & 0 & \frac{1}{3} & 0 & \frac{5}{3}\\
0 & 1 & - \frac{1}{3} & -1 & - \frac{14}{3}\\
0 & 0 & 1 & -6 & -16\\
0 & 0 & 1 & -1 & -1
\end{array}\right]
\]
\end{frame}
\begin{frame}[t]{Ejemplo $4\times4$: paso 9}
Usamos el pivote de la columna 3 para anular la entrada del renglón 1.
\[
R_{1}\leftarrow R_{1}-\frac{1}{3}R_{3}
\]
\[
\left[\begin{array}{cccc|c}
1 & 0 & 0 & 2 & 7\\
0 & 1 & - \frac{1}{3} & -1 & - \frac{14}{3}\\
0 & 0 & 1 & -6 & -16\\
0 & 0 & 1 & -1 & -1
\end{array}\right]
\]
\end{frame}
\begin{frame}[t]{Ejemplo $4\times4$: paso 10}
Usamos el pivote de la columna 3 para anular la entrada del renglón 2.
\[
R_{2}\leftarrow R_{2}+\frac{1}{3}R_{3}
\]
\[
\left[\begin{array}{cccc|c}
1 & 0 & 0 & 2 & 7\\
0 & 1 & 0 & -3 & -10\\
0 & 0 & 1 & -6 & -16\\
0 & 0 & 1 & -1 & -1
\end{array}\right]
\]
\end{frame}
\begin{frame}[t]{Ejemplo $4\times4$: paso 11}
Usamos el pivote de la columna 3 para anular la entrada del renglón 4.
\[
R_{4}\leftarrow R_{4}-R_{3}
\]
\[
\left[\begin{array}{cccc|c}
1 & 0 & 0 & 2 & 7\\
0 & 1 & 0 & -3 & -10\\
0 & 0 & 1 & -6 & -16\\
0 & 0 & 0 & 5 & 15
\end{array}\right]
\]
\end{frame}
\begin{frame}[t]{Ejemplo $4\times4$: paso 12}
Normalizamos el pivote de la columna 4 para convertirlo en 1.
\[
R_{4}\leftarrow \frac{1}{5}R_{4}
\]
\[
\left[\begin{array}{cccc|c}
1 & 0 & 0 & 2 & 7\\
0 & 1 & 0 & -3 & -10\\
0 & 0 & 1 & -6 & -16\\
0 & 0 & 0 & 1 & 3
\end{array}\right]
\]
\end{frame}
\begin{frame}[t]{Ejemplo $4\times4$: paso 13}
Usamos el pivote de la columna 4 para anular la entrada del renglón 1.
\[
R_{1}\leftarrow R_{1}-2R_{4}
\]
\[
\left[\begin{array}{cccc|c}
1 & 0 & 0 & 0 & 1\\
0 & 1 & 0 & -3 & -10\\
0 & 0 & 1 & -6 & -16\\
0 & 0 & 0 & 1 & 3
\end{array}\right]
\]
\end{frame}
\begin{frame}[t]{Ejemplo $4\times4$: paso 14}
Usamos el pivote de la columna 4 para anular la entrada del renglón 2.
\[
R_{2}\leftarrow R_{2}+3R_{4}
\]
\[
\left[\begin{array}{cccc|c}
1 & 0 & 0 & 0 & 1\\
0 & 1 & 0 & 0 & -1\\
0 & 0 & 1 & -6 & -16\\
0 & 0 & 0 & 1 & 3
\end{array}\right]
\]
\end{frame}
\begin{frame}[t]{Ejemplo $4\times4$: paso 15}
Usamos el pivote de la columna 4 para anular la entrada del renglón 3.
\[
R_{3}\leftarrow R_{3}+6R_{4}
\]
\[
\left[\begin{array}{cccc|c}
1 & 0 & 0 & 0 & 1\\
0 & 1 & 0 & 0 & -1\\
0 & 0 & 1 & 0 & 2\\
0 & 0 & 0 & 1 & 3
\end{array}\right]
\]
\end{frame}
\begin{frame}[t]{Ejemplo $4\times4$: conclusión}
La matriz está en forma escalonada reducida. La solución se lee directamente:
\[
\boxed{x_{1}=1,\quad x_{2}=-1,\quad x_{3}=2,\quad x_{4}=3}
\]
\begin{block}{Comprobación}
Sustituimos los valores en el sistema original. Si se recuperan todos los términos independientes, la solución es correcta.
\end{block}
\end{frame}
\subsection{Ejemplo $5\times5$}
\begin{frame}[t]{Ejemplo $5\times5$: planteamiento}
\textbf{Objetivo.} Ejecución exhaustiva de la técnica.
\[
\begin{aligned}
x_{1} + x_{2} - x_{3} + 2 x_{4} + x_{5} &= 8\\
2 x_{1} + 3 x_{2} + x_{3} + x_{4} - 2 x_{5} &= 4\\
- x_{1} + 2 x_{2} + 3 x_{3} - x_{4} + x_{5} &= 1\\
3 x_{1} - x_{2} + 2 x_{3} + 2 x_{4} + x_{5} &= 3\\
2 x_{1} + x_{2} + x_{3} - x_{4} + 3 x_{5} &= 8
\end{aligned}
\]
Matriz aumentada:
\[
\left[\begin{array}{ccccc|c}
1 & 1 & -1 & 2 & 1 & 8\\
2 & 3 & 1 & 1 & -2 & 4\\
-1 & 2 & 3 & -1 & 1 & 1\\
3 & -1 & 2 & 2 & 1 & 3\\
2 & 1 & 1 & -1 & 3 & 8
\end{array}\right]
\]
\end{frame}
\begin{frame}[t]{Ejemplo $5\times5$: paso 1}
Usamos el pivote de la columna 1 para anular la entrada del renglón 2.
\[
R_{2}\leftarrow R_{2}-2R_{1}
\]
\[
\left[\begin{array}{ccccc|c}
1 & 1 & -1 & 2 & 1 & 8\\
0 & 1 & 3 & -3 & -4 & -12\\
-1 & 2 & 3 & -1 & 1 & 1\\
3 & -1 & 2 & 2 & 1 & 3\\
2 & 1 & 1 & -1 & 3 & 8
\end{array}\right]
\]
\end{frame}
\begin{frame}[t]{Ejemplo $5\times5$: paso 2}
Usamos el pivote de la columna 1 para anular la entrada del renglón 3.
\[
R_{3}\leftarrow R_{3}+R_{1}
\]
\[
\left[\begin{array}{ccccc|c}
1 & 1 & -1 & 2 & 1 & 8\\
0 & 1 & 3 & -3 & -4 & -12\\
0 & 3 & 2 & 1 & 2 & 9\\
3 & -1 & 2 & 2 & 1 & 3\\
2 & 1 & 1 & -1 & 3 & 8
\end{array}\right]
\]
\end{frame}
\begin{frame}[t]{Ejemplo $5\times5$: paso 3}
Usamos el pivote de la columna 1 para anular la entrada del renglón 4.
\[
R_{4}\leftarrow R_{4}-3R_{1}
\]
\[
\left[\begin{array}{ccccc|c}
1 & 1 & -1 & 2 & 1 & 8\\
0 & 1 & 3 & -3 & -4 & -12\\
0 & 3 & 2 & 1 & 2 & 9\\
0 & -4 & 5 & -4 & -2 & -21\\
2 & 1 & 1 & -1 & 3 & 8
\end{array}\right]
\]
\end{frame}
\begin{frame}[t]{Ejemplo $5\times5$: paso 4}
Usamos el pivote de la columna 1 para anular la entrada del renglón 5.
\[
R_{5}\leftarrow R_{5}-2R_{1}
\]
\[
\left[\begin{array}{ccccc|c}
1 & 1 & -1 & 2 & 1 & 8\\
0 & 1 & 3 & -3 & -4 & -12\\
0 & 3 & 2 & 1 & 2 & 9\\
0 & -4 & 5 & -4 & -2 & -21\\
0 & -1 & 3 & -5 & 1 & -8
\end{array}\right]
\]
\end{frame}
\begin{frame}[t]{Ejemplo $5\times5$: paso 5}
Usamos el pivote de la columna 2 para anular la entrada del renglón 1.
\[
R_{1}\leftarrow R_{1}-R_{2}
\]
\[
\left[\begin{array}{ccccc|c}
1 & 0 & -4 & 5 & 5 & 20\\
0 & 1 & 3 & -3 & -4 & -12\\
0 & 3 & 2 & 1 & 2 & 9\\
0 & -4 & 5 & -4 & -2 & -21\\
0 & -1 & 3 & -5 & 1 & -8
\end{array}\right]
\]
\end{frame}
\begin{frame}[t]{Ejemplo $5\times5$: paso 6}
Usamos el pivote de la columna 2 para anular la entrada del renglón 3.
\[
R_{3}\leftarrow R_{3}-3R_{2}
\]
\[
\left[\begin{array}{ccccc|c}
1 & 0 & -4 & 5 & 5 & 20\\
0 & 1 & 3 & -3 & -4 & -12\\
0 & 0 & -7 & 10 & 14 & 45\\
0 & -4 & 5 & -4 & -2 & -21\\
0 & -1 & 3 & -5 & 1 & -8
\end{array}\right]
\]
\end{frame}
\begin{frame}[t]{Ejemplo $5\times5$: paso 7}
Usamos el pivote de la columna 2 para anular la entrada del renglón 4.
\[
R_{4}\leftarrow R_{4}+4R_{2}
\]
\[
\left[\begin{array}{ccccc|c}
1 & 0 & -4 & 5 & 5 & 20\\
0 & 1 & 3 & -3 & -4 & -12\\
0 & 0 & -7 & 10 & 14 & 45\\
0 & 0 & 17 & -16 & -18 & -69\\
0 & -1 & 3 & -5 & 1 & -8
\end{array}\right]
\]
\end{frame}
\begin{frame}[t]{Ejemplo $5\times5$: paso 8}
Usamos el pivote de la columna 2 para anular la entrada del renglón 5.
\[
R_{5}\leftarrow R_{5}+R_{2}
\]
\[
\left[\begin{array}{ccccc|c}
1 & 0 & -4 & 5 & 5 & 20\\
0 & 1 & 3 & -3 & -4 & -12\\
0 & 0 & -7 & 10 & 14 & 45\\
0 & 0 & 17 & -16 & -18 & -69\\
0 & 0 & 6 & -8 & -3 & -20
\end{array}\right]
\]
\end{frame}
\begin{frame}[t]{Ejemplo $5\times5$: paso 9}
Normalizamos el pivote de la columna 3 para convertirlo en 1.
\[
R_{3}\leftarrow \frac{1}{-7}R_{3}
\]
\[
\left[\begin{array}{ccccc|c}
1 & 0 & -4 & 5 & 5 & 20\\
0 & 1 & 3 & -3 & -4 & -12\\
0 & 0 & 1 & - \frac{10}{7} & -2 & - \frac{45}{7}\\
0 & 0 & 17 & -16 & -18 & -69\\
0 & 0 & 6 & -8 & -3 & -20
\end{array}\right]
\]
\end{frame}
\begin{frame}[t]{Ejemplo $5\times5$: paso 10}
Usamos el pivote de la columna 3 para anular la entrada del renglón 1.
\[
R_{1}\leftarrow R_{1}+4R_{3}
\]
\[
\left[\begin{array}{ccccc|c}
1 & 0 & 0 & - \frac{5}{7} & -3 & - \frac{40}{7}\\
0 & 1 & 3 & -3 & -4 & -12\\
0 & 0 & 1 & - \frac{10}{7} & -2 & - \frac{45}{7}\\
0 & 0 & 17 & -16 & -18 & -69\\
0 & 0 & 6 & -8 & -3 & -20
\end{array}\right]
\]
\end{frame}
\begin{frame}[t]{Ejemplo $5\times5$: paso 11}
Usamos el pivote de la columna 3 para anular la entrada del renglón 2.
\[
R_{2}\leftarrow R_{2}-3R_{3}
\]
\[
\left[\begin{array}{ccccc|c}
1 & 0 & 0 & - \frac{5}{7} & -3 & - \frac{40}{7}\\
0 & 1 & 0 & \frac{9}{7} & 2 & \frac{51}{7}\\
0 & 0 & 1 & - \frac{10}{7} & -2 & - \frac{45}{7}\\
0 & 0 & 17 & -16 & -18 & -69\\
0 & 0 & 6 & -8 & -3 & -20
\end{array}\right]
\]
\end{frame}
\begin{frame}[t]{Ejemplo $5\times5$: paso 12}
Usamos el pivote de la columna 3 para anular la entrada del renglón 4.
\[
R_{4}\leftarrow R_{4}-17R_{3}
\]
\[
\left[\begin{array}{ccccc|c}
1 & 0 & 0 & - \frac{5}{7} & -3 & - \frac{40}{7}\\
0 & 1 & 0 & \frac{9}{7} & 2 & \frac{51}{7}\\
0 & 0 & 1 & - \frac{10}{7} & -2 & - \frac{45}{7}\\
0 & 0 & 0 & \frac{58}{7} & 16 & \frac{282}{7}\\
0 & 0 & 6 & -8 & -3 & -20
\end{array}\right]
\]
\end{frame}
\begin{frame}[t]{Ejemplo $5\times5$: paso 13}
Usamos el pivote de la columna 3 para anular la entrada del renglón 5.
\[
R_{5}\leftarrow R_{5}-6R_{3}
\]
\[
\left[\begin{array}{ccccc|c}
1 & 0 & 0 & - \frac{5}{7} & -3 & - \frac{40}{7}\\
0 & 1 & 0 & \frac{9}{7} & 2 & \frac{51}{7}\\
0 & 0 & 1 & - \frac{10}{7} & -2 & - \frac{45}{7}\\
0 & 0 & 0 & \frac{58}{7} & 16 & \frac{282}{7}\\
0 & 0 & 0 & \frac{4}{7} & 9 & \frac{130}{7}
\end{array}\right]
\]
\end{frame}
\begin{frame}[t]{Ejemplo $5\times5$: paso 14}
Normalizamos el pivote de la columna 4 para convertirlo en 1.
\[
R_{4}\leftarrow \frac{1}{\frac{58}{7}}R_{4}
\]
\[
\left[\begin{array}{ccccc|c}
1 & 0 & 0 & - \frac{5}{7} & -3 & - \frac{40}{7}\\
0 & 1 & 0 & \frac{9}{7} & 2 & \frac{51}{7}\\
0 & 0 & 1 & - \frac{10}{7} & -2 & - \frac{45}{7}\\
0 & 0 & 0 & 1 & \frac{56}{29} & \frac{141}{29}\\
0 & 0 & 0 & \frac{4}{7} & 9 & \frac{130}{7}
\end{array}\right]
\]
\end{frame}
\begin{frame}[t]{Ejemplo $5\times5$: paso 15}
Usamos el pivote de la columna 4 para anular la entrada del renglón 1.
\[
R_{1}\leftarrow R_{1}+\frac{5}{7}R_{4}
\]
\[
\left[\begin{array}{ccccc|c}
1 & 0 & 0 & 0 & - \frac{47}{29} & - \frac{65}{29}\\
0 & 1 & 0 & \frac{9}{7} & 2 & \frac{51}{7}\\
0 & 0 & 1 & - \frac{10}{7} & -2 & - \frac{45}{7}\\
0 & 0 & 0 & 1 & \frac{56}{29} & \frac{141}{29}\\
0 & 0 & 0 & \frac{4}{7} & 9 & \frac{130}{7}
\end{array}\right]
\]
\end{frame}
\begin{frame}[t]{Ejemplo $5\times5$: paso 16}
Usamos el pivote de la columna 4 para anular la entrada del renglón 2.
\[
R_{2}\leftarrow R_{2}-\frac{9}{7}R_{4}
\]
\[
\left[\begin{array}{ccccc|c}
1 & 0 & 0 & 0 & - \frac{47}{29} & - \frac{65}{29}\\
0 & 1 & 0 & 0 & - \frac{14}{29} & \frac{30}{29}\\
0 & 0 & 1 & - \frac{10}{7} & -2 & - \frac{45}{7}\\
0 & 0 & 0 & 1 & \frac{56}{29} & \frac{141}{29}\\
0 & 0 & 0 & \frac{4}{7} & 9 & \frac{130}{7}
\end{array}\right]
\]
\end{frame}
\begin{frame}[t]{Ejemplo $5\times5$: paso 17}
Usamos el pivote de la columna 4 para anular la entrada del renglón 3.
\[
R_{3}\leftarrow R_{3}+\frac{10}{7}R_{4}
\]
\[
\left[\begin{array}{ccccc|c}
1 & 0 & 0 & 0 & - \frac{47}{29} & - \frac{65}{29}\\
0 & 1 & 0 & 0 & - \frac{14}{29} & \frac{30}{29}\\
0 & 0 & 1 & 0 & \frac{22}{29} & \frac{15}{29}\\
0 & 0 & 0 & 1 & \frac{56}{29} & \frac{141}{29}\\
0 & 0 & 0 & \frac{4}{7} & 9 & \frac{130}{7}
\end{array}\right]
\]
\end{frame}
\begin{frame}[t]{Ejemplo $5\times5$: paso 18}
Usamos el pivote de la columna 4 para anular la entrada del renglón 5.
\[
R_{5}\leftarrow R_{5}-\frac{4}{7}R_{4}
\]
\[
\left[\begin{array}{ccccc|c}
1 & 0 & 0 & 0 & - \frac{47}{29} & - \frac{65}{29}\\
0 & 1 & 0 & 0 & - \frac{14}{29} & \frac{30}{29}\\
0 & 0 & 1 & 0 & \frac{22}{29} & \frac{15}{29}\\
0 & 0 & 0 & 1 & \frac{56}{29} & \frac{141}{29}\\
0 & 0 & 0 & 0 & \frac{229}{29} & \frac{458}{29}
\end{array}\right]
\]
\end{frame}
\begin{frame}[t]{Ejemplo $5\times5$: paso 19}
Normalizamos el pivote de la columna 5 para convertirlo en 1.
\[
R_{5}\leftarrow \frac{1}{\frac{229}{29}}R_{5}
\]
\[
\left[\begin{array}{ccccc|c}
1 & 0 & 0 & 0 & - \frac{47}{29} & - \frac{65}{29}\\
0 & 1 & 0 & 0 & - \frac{14}{29} & \frac{30}{29}\\
0 & 0 & 1 & 0 & \frac{22}{29} & \frac{15}{29}\\
0 & 0 & 0 & 1 & \frac{56}{29} & \frac{141}{29}\\
0 & 0 & 0 & 0 & 1 & 2
\end{array}\right]
\]
\end{frame}
\begin{frame}[t]{Ejemplo $5\times5$: paso 20}
Usamos el pivote de la columna 5 para anular la entrada del renglón 1.
\[
R_{1}\leftarrow R_{1}+\frac{47}{29}R_{5}
\]
\[
\left[\begin{array}{ccccc|c}
1 & 0 & 0 & 0 & 0 & 1\\
0 & 1 & 0 & 0 & - \frac{14}{29} & \frac{30}{29}\\
0 & 0 & 1 & 0 & \frac{22}{29} & \frac{15}{29}\\
0 & 0 & 0 & 1 & \frac{56}{29} & \frac{141}{29}\\
0 & 0 & 0 & 0 & 1 & 2
\end{array}\right]
\]
\end{frame}
\begin{frame}[t]{Ejemplo $5\times5$: paso 21}
Usamos el pivote de la columna 5 para anular la entrada del renglón 2.
\[
R_{2}\leftarrow R_{2}+\frac{14}{29}R_{5}
\]
\[
\left[\begin{array}{ccccc|c}
1 & 0 & 0 & 0 & 0 & 1\\
0 & 1 & 0 & 0 & 0 & 2\\
0 & 0 & 1 & 0 & \frac{22}{29} & \frac{15}{29}\\
0 & 0 & 0 & 1 & \frac{56}{29} & \frac{141}{29}\\
0 & 0 & 0 & 0 & 1 & 2
\end{array}\right]
\]
\end{frame}
\begin{frame}[t]{Ejemplo $5\times5$: paso 22}
Usamos el pivote de la columna 5 para anular la entrada del renglón 3.
\[
R_{3}\leftarrow R_{3}-\frac{22}{29}R_{5}
\]
\[
\left[\begin{array}{ccccc|c}
1 & 0 & 0 & 0 & 0 & 1\\
0 & 1 & 0 & 0 & 0 & 2\\
0 & 0 & 1 & 0 & 0 & -1\\
0 & 0 & 0 & 1 & \frac{56}{29} & \frac{141}{29}\\
0 & 0 & 0 & 0 & 1 & 2
\end{array}\right]
\]
\end{frame}
\begin{frame}[t]{Ejemplo $5\times5$: paso 23}
Usamos el pivote de la columna 5 para anular la entrada del renglón 4.
\[
R_{4}\leftarrow R_{4}-\frac{56}{29}R_{5}
\]
\[
\left[\begin{array}{ccccc|c}
1 & 0 & 0 & 0 & 0 & 1\\
0 & 1 & 0 & 0 & 0 & 2\\
0 & 0 & 1 & 0 & 0 & -1\\
0 & 0 & 0 & 1 & 0 & 1\\
0 & 0 & 0 & 0 & 1 & 2
\end{array}\right]
\]
\end{frame}
\begin{frame}[t]{Ejemplo $5\times5$: conclusión}
La matriz está en forma escalonada reducida. La solución se lee directamente:
\[
\boxed{x_{1}=1,\quad x_{2}=2,\quad x_{3}=-1,\quad x_{4}=1,\quad x_{5}=2}
\]
\begin{block}{Comprobación}
Sustituimos los valores en el sistema original. Si se recuperan todos los términos independientes, la solución es correcta.
\end{block}
\end{frame}

\begin{frame}[t]{Gauss--Jordan: infinitas soluciones}
Una RREF posible es
\[
\left[\begin{array}{cccc|c}
1&0&-3&3&2\\
0&1&1&-1&1\\
0&0&0&0&0\\
0&0&0&0&0
\end{array}\right].
\]
Las variables \(x_3=s\) y \(x_4=t\) son libres:
\[
\boxed{(x_1,x_2,x_3,x_4)=(2+3s-3t,\,1-s+t,\,s,\,t)}.
\]
\end{frame}

\begin{frame}[t]{Gauss--Jordan: incompatibilidad}
Si durante la reducción aparece
\[
\left[\begin{array}{ccc|c}
*&*&*&*\\
*&*&*&*\\
0&0&0&1
\end{array}\right],
\]
el último renglón representa \(0=1\). Por tanto,
\[
\boxed{\text{el sistema no tiene solución}.}
\]
\end{frame}

\subsection{Gauss frente a Gauss--Jordan}
\begin{frame}[t]{Gauss frente a Gauss--Jordan}
Gauss termina en forma escalonada y utiliza sustitución hacia atrás. Gauss--Jordan continúa hasta RREF.
En sistemas grandes, Gauss suele requerir menos operaciones; Gauss--Jordan hace especialmente visibles
las variables básicas, las variables libres, el rango y la consistencia.
\end{frame}

\section{Sistemas con parámetros}
\begin{frame}[allowframebreaks,t]{Sistemas con parámetros}
Los parámetros permiten estudiar familias completas de sistemas.

\par\medskip{\large\bfseries Parámetros como variables libres}\par\smallskip

Consideremos
\[
\begin{cases}
x+y+z=4,\\
2x+2y+2z=8.
\end{cases}
\]
Después de reducir,
\[
x+y+z=4.
\]
Tomando
\[
y=s,\qquad z=t,
\]
obtenemos
\[
x=4-s-t.
\]
Aquí $s$ y $t$ son parámetros que describen todas las soluciones.

\par\medskip{\large\bfseries Parámetros en los coeficientes}\par\smallskip

Consideremos
\[
\begin{cases}
x+y=2,\\
2x+ay=4.
\end{cases}
\]
La matriz aumentada es
\[
\left[
\begin{array}{cc|c}
1&1&2\\
2&a&4
\end{array}
\right].
\]
Aplicando
\[
R_2\leftarrow R_2-2R_1
\]
obtenemos
\[
\left[
\begin{array}{cc|c}
1&1&2\\
0&a-2&0
\end{array}
\right].
\]

Si $a\neq2$, entonces
\[
(a-2)y=0\quad\Longrightarrow\quad y=0,
\]
y $x=2$. Existe una única solución.

Si $a=2$, el segundo renglón se vuelve nulo:
\[
[0\quad0\mid0],
\]
y queda una ecuación con dos incógnitas. Existen infinitas soluciones.

Este ejemplo muestra por qué, al trabajar con parámetros, no debemos dividir por una expresión como
$a-2$ sin estudiar antes el caso en que puede ser cero.

\par\medskip{\large\bfseries Un parámetro que produce incompatibilidad}\par\smallskip

\[
\begin{cases}
x+y=2,\\
2x+2y=k.
\end{cases}
\]
Reduciendo:
\[
\left[
\begin{array}{cc|c}
1&1&2\\
0&0&k-4
\end{array}
\right].
\]
Por tanto:
\[
\begin{cases}
k=4 &\Longrightarrow \text{infinitas soluciones},\\
k\neq4 &\Longrightarrow \text{ninguna solución}.
\end{cases}
\]
\end{frame}

\section{Rango de una matriz}
\begin{frame}[allowframebreaks,t]{Rango de una matriz}
\begin{definicion}
El \textbf{rango} de una matriz $A$, denotado por $\rank(A)$, es el número de pivotes de una forma
escalonada de $A$.
\end{definicion}

Equivalentemente, es el número de renglones no nulos de su forma escalonada.

\begin{ejemplo}
Si
\[
A\sim
\begin{pmatrix}
1&2&0&3\\
0&1&4&2\\
0&0&0&0
\end{pmatrix},
\]
entonces
\[
\rank(A)=2.
\]
\end{ejemplo}
\end{frame}

\section{Rango y consistencia}
\begin{frame}[allowframebreaks,t]{Rango y consistencia}
Para el sistema
\[
A\mathbf{x}=\mathbf b,
\]
debemos comparar la matriz de coeficientes $A$ con la matriz aumentada $[A\mid\mathbf b]$.

\begin{teorema}[Criterio de consistencia]
El sistema $A\mathbf{x}=\mathbf b$ es compatible si y sólo si
\[
\rank(A)=\rank([A\mid\mathbf b]).
\]
\end{teorema}

\begin{observacion}
Si al reducir la matriz aumentada aparece un renglón
\[
[0\ \cdots\ 0\mid c],\qquad c\neq0,
\]
entonces la última columna introduce un pivote que no existe en $A$. Por ello
\[
\rank(A)<\rank([A\mid\mathbf b]),
\]
y el sistema es incompatible.
\end{observacion}

Si el sistema es compatible y tiene $n$ incógnitas:
\[
\begin{cases}
\rank(A)=n &\Longrightarrow \text{solución única},\\
\rank(A)<n &\Longrightarrow \text{infinitas soluciones}.
\end{cases}
\]
\end{frame}

\section{Solución mediante matriz inversa}
\begin{frame}[allowframebreaks,t]{Solución mediante matriz inversa}
Supongamos que $A$ es una matriz cuadrada e invertible. Si
\[
A\mathbf{x}=\mathbf b,
\]
multiplicamos por $A^{-1}$:
\[
A^{-1}A\mathbf{x}=A^{-1}\mathbf b.
\]
Como
\[
A^{-1}A=I,
\]
obtenemos
\[
\boxed{\mathbf{x}=A^{-1}\mathbf b.}
\]

\begin{ejemplo}
Sea
\[
A=
\begin{pmatrix}
2&1\\
1&1
\end{pmatrix},
\qquad
\mathbf b=
\begin{pmatrix}5\\3\end{pmatrix}.
\]
Como
\[
A^{-1}=
\begin{pmatrix}
1&-1\\
-1&2
\end{pmatrix},
\]
se tiene
\[
\mathbf x=A^{-1}\mathbf b
=
\begin{pmatrix}
1&-1\\
-1&2
\end{pmatrix}
\begin{pmatrix}5\\3\end{pmatrix}
=
\begin{pmatrix}2\\1\end{pmatrix}.
\]
\end{ejemplo}

\begin{observacion}
La fórmula $A^{-1}\mathbf b$ sólo tiene sentido cuando $A$ es cuadrada e invertible. Además, calcular
explícitamente la inversa no suele ser el procedimiento más eficiente para resolver un único sistema.
El método es conceptualmente importante porque relaciona invertibilidad con existencia y unicidad.
\end{observacion}
\end{frame}

\section{Regla de Cramer}
\begin{frame}[allowframebreaks,t]{Regla de Cramer}
La regla de Cramer proporciona una fórmula para sistemas cuadrados cuya matriz de coeficientes tiene
determinante distinto de cero.

Sea
\[
A\mathbf x=\mathbf b,\qquad \det(A)\neq0.
\]
Para cada $i$, sea $A_i$ la matriz obtenida al sustituir la columna $i$ de $A$ por $\mathbf b$. Entonces
\[
\boxed{x_i=\frac{\det(A_i)}{\det(A)}}.
\]

\end{frame}
\begin{frame}[t,shrink=18]{Regla de Cramer: ejemplo}
\begin{ejemplo}
Resolvamos
\[
\begin{cases}
2x+y=5,\\
x-y=1.
\end{cases}
\]
Tenemos
\[
A=
\begin{pmatrix}2&1\\1&-1\end{pmatrix},
\qquad
\det(A)=-3.
\]
Para $x$,
\[
A_1=
\begin{pmatrix}5&1\\1&-1\end{pmatrix},
\qquad
\det(A_1)=-6,
\]
por lo que
\[
x=\frac{-6}{-3}=2.
\]
Para $y$,
\[
A_2=
\begin{pmatrix}2&5\\1&1\end{pmatrix},
\qquad
\det(A_2)=-3,
\]
y
\[
y=\frac{-3}{-3}=1.
\]
\end{ejemplo}

\begin{observacion}
Cramer requiere $\det(A)\neq0$. Es útil teóricamente y para sistemas pequeños, pero para sistemas
grandes la eliminación gaussiana suele ser más eficiente.
\end{observacion}
\end{frame}

\section{Cómo elegir un método}
\begin{frame}[allowframebreaks,t]{Cómo elegir un método}
No existe una única técnica que deba aplicarse mecánicamente a todos los sistemas.

\begin{center}
\begin{tabular}{p{0.35\textwidth}p{0.52\textwidth}}
\toprule
\textbf{Situación} & \textbf{Método o idea conveniente}\\
\midrule
Sistema general & Eliminación de Gauss\\
Se desea RREF, variables libres o estructura completa & Gauss--Jordan\\
Matriz cuadrada e inversa ya conocida & $\mathbf x=A^{-1}\mathbf b$\\
Sistema cuadrado pequeño, $\det A\neq0$ & Regla de Cramer\\
Coeficientes con parámetros & Gauss, separando cuidadosamente casos\\
Interés en consistencia & Comparación de rangos\\
\bottomrule
\end{tabular}
\end{center}
\end{frame}

\section{Aplicaciones}
\begin{frame}[allowframebreaks,t]{Aplicaciones}
\par\medskip{\large\bfseries Mezclas}\par\smallskip

Una solución de $10$ litros debe contener $30\%$ de cierta sustancia. Se dispone de una solución al
$20\%$ y otra al $50\%$. Si $x$ y $y$ son los litros utilizados de cada una,
\[
x+y=10.
\]
La cantidad de sustancia pura debe satisfacer
\[
0.20x+0.50y=0.30(10)=3.
\]
Así,
\[
\begin{cases}
x+y=10,\\
0.2x+0.5y=3.
\end{cases}
\]
Multiplicando la segunda ecuación por $10$:
\[
2x+5y=30.
\]
De $x=10-y$:
\[
2(10-y)+5y=30,
\]
por lo que $3y=10$ y
\[
y=\frac{10}{3},\qquad x=\frac{20}{3}.
\]

\par\medskip{\large\bfseries Circuitos}\par\smallskip

En muchos circuitos, las leyes de Kirchhoff producen sistemas lineales en las corrientes desconocidas.
Por ejemplo,
\[
\begin{cases}
I_1-I_2-I_3=0,\\
4I_1+2I_2=12,\\
2I_2-3I_3=0.
\end{cases}
\]
La primera ecuación puede representar conservación de corriente en un nodo; las restantes, balances de
voltaje en mallas. Una vez formulado el modelo, el problema se convierte en resolver un sistema lineal.

\par\medskip{\large\bfseries Equilibrio de fuerzas}\par\smallskip

Si varias fuerzas actúan sobre una estructura en equilibrio, la suma de componentes debe ser cero:
\[
\sum F_x=0,\qquad \sum F_y=0.
\]
Cuando las magnitudes desconocidas aparecen linealmente, estas ecuaciones forman un sistema lineal.
\end{frame}

\section{Errores frecuentes}
\begin{frame}[allowframebreaks,t]{Errores frecuentes}
\begin{enumerate}
\item Realizar una operación en el lado izquierdo de una ecuación sin hacerla también en el término
independiente. La matriz aumentada evita parcialmente este error si se opera el renglón completo.
\item Dividir por una expresión que contiene un parámetro sin estudiar primero cuándo esa expresión
puede ser cero.
\item Confundir un renglón nulo,
\[
[0\ 0\ 0\mid0],
\]
con una contradicción. El renglón nulo no impone una nueva condición.
\item No reconocer una contradicción:
\[
[0\ 0\ 0\mid c],\qquad c\neq0.
\]
\item Asignar arbitrariamente el valor cero a las variables libres. Eso produce una solución particular,
pero no describe el conjunto completo de soluciones.
\item Usar la matriz inversa o Cramer cuando sus hipótesis no se cumplen.
\item Detener Gauss demasiado pronto: para estar en forma escalonada deben existir ceros debajo de
todos los pivotes.
\end{enumerate}
\end{frame}

\section{Ejemplos integradores}
\begin{frame}[allowframebreaks,t]{Ejemplos integradores}
\par\medskip{\large\bfseries Ejemplo A: clasificar y resolver}\par\smallskip

\[
\begin{cases}
x+2y-z=1,\\
2x+4y-2z=2,\\
x+y+z=3.
\end{cases}
\]
La matriz aumentada es
\[
\left[
\begin{array}{ccc|c}
1&2&-1&1\\
2&4&-2&2\\
1&1&1&3
\end{array}
\right].
\]
Aplicamos
\[
R_2\leftarrow R_2-2R_1,\qquad
R_3\leftarrow R_3-R_1,
\]
y obtenemos
\[
\left[
\begin{array}{ccc|c}
1&2&-1&1\\
0&0&0&0\\
0&-1&2&2
\end{array}
\right].
\]
Intercambiando los dos últimos renglones:
\[
\left[
\begin{array}{ccc|c}
1&2&-1&1\\
0&-1&2&2\\
0&0&0&0
\end{array}
\right].
\]
Hay dos pivotes y tres incógnitas. Por tanto existe una variable libre. Tomemos $z=t$. La segunda
ecuación da
\[
-y+2t=2\quad\Longrightarrow\quad y=2t-2.
\]
La primera:
\[
x+2(2t-2)-t=1,
\]
de donde
\[
x=5-3t.
\]
Así,
\[
\boxed{(x,y,z)=(5-3t,\,2t-2,\,t),\quad t\in\R.}
\]
El sistema es compatible indeterminado.

\par\medskip{\large\bfseries Ejemplo B: parámetro y clasificación}\par\smallskip

\[
\begin{cases}
x+y+z=1,\\
2x+2y+2z=2,\\
x+y+z=k.
\end{cases}
\]
Restando la primera ecuación de la tercera obtenemos
\[
0=k-1.
\]
Por tanto:
\[
k=1\Longrightarrow\text{infinitas soluciones},
\]
mientras que
\[
k\neq1\Longrightarrow\text{sistema incompatible}.
\]
\end{frame}

\section{Banco de ejercicios para Gauss--Jordan}
\begin{frame}[t]{Cómo utilizar el banco}
El banco está pensado para elegir ejercicios durante la clase, no para resolverlos todos en una sesión.
\begin{itemize}
\item \(2\times2\): mecánica básica.
\item \(3\times3\): algoritmo completo.
\item \(4\times4\): organización por pivotes.
\item \(5\times5\): ejecución exhaustiva.
\item Casos especiales: variables libres e incompatibilidad.
\end{itemize}
\end{frame}

\begin{frame}[t]{Banco Gauss--Jordan: ejercicio 1 (2$\times$2)}
\textbf{Nivel:} Básico.\\
\textbf{Objetivo:} Mecánica de operaciones elementales.
\[
\begin{aligned}
3 x_{1} - 2 x_{2} &= 7\\
2 x_{1} - x_{2} &= 5
\end{aligned}
\]
\end{frame}
\begin{frame}[t]{Banco Gauss--Jordan: ejercicio 2 (2$\times$2)}
\textbf{Nivel:} Básico.\\
\textbf{Objetivo:} Mecánica de operaciones elementales.
\[
\begin{aligned}
- 3 x_{1} + 4 x_{2} &= 10\\
- 4 x_{1} &= -8
\end{aligned}
\]
\end{frame}
\begin{frame}[t]{Banco Gauss--Jordan: ejercicio 3 (2$\times$2)}
\textbf{Nivel:} Básico.\\
\textbf{Objetivo:} Mecánica de operaciones elementales.
\[
\begin{aligned}
3 x_{1} + 4 x_{2} &= 25\\
3 x_{1} - 2 x_{2} &= 1
\end{aligned}
\]
\end{frame}
\begin{frame}[t]{Banco Gauss--Jordan: ejercicio 4 (2$\times$2)}
\textbf{Nivel:} Básico.\\
\textbf{Objetivo:} Mecánica de operaciones elementales.
\[
\begin{aligned}
3 x_{1} + 2 x_{2} &= -11\\
- 2 x_{1} - 2 x_{2} &= 8
\end{aligned}
\]
\end{frame}
\begin{frame}[t]{Banco Gauss--Jordan: ejercicio 5 (2$\times$2)}
\textbf{Nivel:} Básico.\\
\textbf{Objetivo:} Mecánica de operaciones elementales.
\[
\begin{aligned}
3 x_{1} - 2 x_{2} &= -1\\
- 4 x_{1} - 2 x_{2} &= -8
\end{aligned}
\]
\end{frame}
\begin{frame}[t]{Banco Gauss--Jordan: ejercicio 6 (3$\times$3)}
\textbf{Nivel:} Básico/Intermedio.\\
\textbf{Objetivo:} Algoritmo completo y lectura de la RREF.
\[
\begin{aligned}
- x_{2} + 2 x_{3} &= 0\\
3 x_{1} - x_{2} - 4 x_{3} &= 0\\
x_{1} - 4 x_{2} + 2 x_{3} &= -8
\end{aligned}
\]
\end{frame}
\begin{frame}[t]{Banco Gauss--Jordan: ejercicio 7 (3$\times$3)}
\textbf{Nivel:} Básico/Intermedio.\\
\textbf{Objetivo:} Algoritmo completo y lectura de la RREF.
\[
\begin{aligned}
2 x_{1} + 3 x_{2} + 3 x_{3} &= 8\\
- x_{1} - 4 x_{2} - 2 x_{3} &= -5\\
- 2 x_{1} - 2 x_{2} - 2 x_{3} &= -6
\end{aligned}
\]
\end{frame}
\begin{frame}[t]{Banco Gauss--Jordan: ejercicio 8 (3$\times$3)}
\textbf{Nivel:} Básico/Intermedio.\\
\textbf{Objetivo:} Algoritmo completo y lectura de la RREF.
\[
\begin{aligned}
- 3 x_{1} - 4 x_{2} - 2 x_{3} &= 6\\
x_{2} - x_{3} &= -6\\
- x_{1} - 2 x_{2} - 2 x_{3} &= -2
\end{aligned}
\]
\end{frame}
\begin{frame}[t]{Banco Gauss--Jordan: ejercicio 9 (3$\times$3)}
\textbf{Nivel:} Básico/Intermedio.\\
\textbf{Objetivo:} Algoritmo completo y lectura de la RREF.
\[
\begin{aligned}
- 2 x_{1} - 4 x_{2} - 3 x_{3} &= -17\\
2 x_{1} - 2 x_{2} &= -4\\
4 x_{1} + 3 x_{2} + 4 x_{3} &= 17
\end{aligned}
\]
\end{frame}
\begin{frame}[t]{Banco Gauss--Jordan: ejercicio 10 (3$\times$3)}
\textbf{Nivel:} Básico/Intermedio.\\
\textbf{Objetivo:} Algoritmo completo y lectura de la RREF.
\[
\begin{aligned}
- 4 x_{1} + 4 x_{3} &= 4\\
- x_{1} - 3 x_{2} + 3 x_{3} &= 8\\
2 x_{3} &= -2
\end{aligned}
\]
\end{frame}
\begin{frame}[t]{Banco Gauss--Jordan: ejercicio 11 (3$\times$3)}
\textbf{Nivel:} Básico/Intermedio.\\
\textbf{Objetivo:} Algoritmo completo y lectura de la RREF.
\[
\begin{aligned}
- 2 x_{1} + 3 x_{2} - 3 x_{3} &= -11\\
2 x_{1} - 3 x_{2} - 3 x_{3} &= -1\\
- 4 x_{1} - 4 x_{2} - x_{3} &= -2
\end{aligned}
\]
\end{frame}
\begin{frame}[t]{Banco Gauss--Jordan: ejercicio 12 (3$\times$3)}
\textbf{Nivel:} Básico/Intermedio.\\
\textbf{Objetivo:} Algoritmo completo y lectura de la RREF.
\[
\begin{aligned}
- 3 x_{1} + x_{2} &= -8\\
3 x_{1} - 2 x_{2} + 2 x_{3} &= 16\\
2 x_{2} - x_{3} &= -7
\end{aligned}
\]
\end{frame}
\begin{frame}[t]{Banco Gauss--Jordan: ejercicio 13 (3$\times$3)}
\textbf{Nivel:} Básico/Intermedio.\\
\textbf{Objetivo:} Algoritmo completo y lectura de la RREF.
\[
\begin{aligned}
x_{1} - 2 x_{2} - 4 x_{3} &= 2\\
- 4 x_{1} + 2 x_{2} + x_{3} &= 1\\
- 2 x_{1} - x_{2} + 4 x_{3} &= -5
\end{aligned}
\]
\end{frame}
\begin{frame}[t]{Banco Gauss--Jordan: ejercicio 14 (4$\times$4)}
\textbf{Nivel:} Intermedio.\\
\textbf{Objetivo:} Organización por pivotes en sistemas medianos.
\[
\begin{aligned}
- x_{1} + 3 x_{2} + 2 x_{3} - x_{4} &= 5\\
x_{1} + x_{2} - 2 x_{3} + 2 x_{4} &= -2\\
- 3 x_{1} - 2 x_{2} - 2 x_{3} + 2 x_{4} &= -1\\
- x_{1} - 2 x_{2} + 3 x_{3} - 3 x_{4} &= 2
\end{aligned}
\]
\end{frame}
\begin{frame}[t]{Banco Gauss--Jordan: ejercicio 15 (4$\times$4)}
\textbf{Nivel:} Intermedio.\\
\textbf{Objetivo:} Organización por pivotes en sistemas medianos.
\[
\begin{aligned}
- 3 x_{1} - 2 x_{2} + 2 x_{3} &= 0\\
- 2 x_{1} + 3 x_{2} - 3 x_{3} + 3 x_{4} &= 25\\
- 3 x_{1} + x_{3} - 2 x_{4} &= -5\\
2 x_{1} - 2 x_{2} - 3 x_{3} &= 5
\end{aligned}
\]
\end{frame}
\begin{frame}[t]{Banco Gauss--Jordan: ejercicio 16 (4$\times$4)}
\textbf{Nivel:} Intermedio.\\
\textbf{Objetivo:} Organización por pivotes en sistemas medianos.
\[
\begin{aligned}
x_{1} - x_{2} - 3 x_{4} &= -5\\
- 3 x_{1} - 3 x_{2} + 3 x_{3} + 2 x_{4} &= 7\\
- 3 x_{2} - 3 x_{3} - x_{4} &= -8\\
- 2 x_{1} + 3 x_{2} - 3 x_{3} + x_{4} &= -6
\end{aligned}
\]
\end{frame}
\begin{frame}[t]{Banco Gauss--Jordan: ejercicio 17 (4$\times$4)}
\textbf{Nivel:} Intermedio.\\
\textbf{Objetivo:} Organización por pivotes en sistemas medianos.
\[
\begin{aligned}
- 3 x_{1} - 2 x_{2} - 2 x_{3} &= 9\\
x_{1} + x_{2} + 3 x_{3} + x_{4} &= -9\\
3 x_{3} + x_{4} &= -9\\
x_{1} - 2 x_{3} + x_{4} &= 3
\end{aligned}
\]
\end{frame}
\begin{frame}[t]{Banco Gauss--Jordan: ejercicio 18 (4$\times$4)}
\textbf{Nivel:} Intermedio.\\
\textbf{Objetivo:} Organización por pivotes en sistemas medianos.
\[
\begin{aligned}
2 x_{1} - 2 x_{2} - x_{3} &= -10\\
2 x_{1} - x_{2} - x_{3} + x_{4} &= -7\\
x_{1} - 3 x_{3} - 2 x_{4} &= -4\\
3 x_{2} - 3 x_{3} + 3 x_{4} &= 3
\end{aligned}
\]
\end{frame}
\begin{frame}[t]{Banco Gauss--Jordan: ejercicio 19 (4$\times$4)}
\textbf{Nivel:} Intermedio.\\
\textbf{Objetivo:} Organización por pivotes en sistemas medianos.
\[
\begin{aligned}
- x_{1} + 2 x_{2} - x_{3} - 2 x_{4} &= -14\\
3 x_{1} + 2 x_{3} - x_{4} &= 13\\
3 x_{1} - 3 x_{2} - x_{3} - 2 x_{4} &= 8\\
- 3 x_{1} + x_{2} - x_{4} &= -13
\end{aligned}
\]
\end{frame}
\begin{frame}[t]{Banco Gauss--Jordan: ejercicio 20 (4$\times$4)}
\textbf{Nivel:} Intermedio.\\
\textbf{Objetivo:} Organización por pivotes en sistemas medianos.
\[
\begin{aligned}
3 x_{2} + 3 x_{3} - 2 x_{4} &= 9\\
2 x_{1} - x_{2} - 2 x_{3} + x_{4} &= -3\\
3 x_{1} + 2 x_{2} + 3 x_{4} &= -7\\
x_{2} - x_{3} &= -5
\end{aligned}
\]
\end{frame}
\begin{frame}[t]{Banco Gauss--Jordan: ejercicio 21 (5$\times$5)}
\textbf{Nivel:} Avanzado.\\
\textbf{Objetivo:} Ejecución exhaustiva del algoritmo.
\[
\begin{aligned}
x_{1} - x_{2} + 2 x_{3} - 2 x_{5} &= 6\\
- 3 x_{1} + 3 x_{2} - x_{3} + x_{4} &= 14\\
- x_{1} - x_{2} + x_{3} - 3 x_{4} + 3 x_{5} &= -3\\
2 x_{1} - 3 x_{2} - 3 x_{3} - 3 x_{4} &= -21\\
3 x_{1} + x_{2} - 2 x_{3} + 2 x_{4} + x_{5} &= -17
\end{aligned}
\]
\end{frame}
\begin{frame}[t]{Banco Gauss--Jordan: ejercicio 22 (5$\times$5)}
\textbf{Nivel:} Avanzado.\\
\textbf{Objetivo:} Ejecución exhaustiva del algoritmo.
\[
\begin{aligned}
3 x_{1} - 3 x_{2} + x_{3} - 2 x_{4} - x_{5} &= 16\\
- x_{1} - 3 x_{2} + 2 x_{3} + 3 x_{4} - 2 x_{5} &= 1\\
- 2 x_{1} - 2 x_{2} + 3 x_{3} + 3 x_{4} + 2 x_{5} &= -4\\
- 2 x_{1} + x_{2} + 2 x_{3} + x_{4} - 3 x_{5} &= 0\\
2 x_{1} + x_{2} + 2 x_{3} + x_{4} - 3 x_{5} &= 8
\end{aligned}
\]
\end{frame}
\begin{frame}[t]{Banco Gauss--Jordan: ejercicio 23 (5$\times$5)}
\textbf{Nivel:} Avanzado.\\
\textbf{Objetivo:} Ejecución exhaustiva del algoritmo.
\[
\begin{aligned}
- 3 x_{1} - x_{2} - 2 x_{3} + 2 x_{4} + x_{5} &= 2\\
3 x_{1} + 3 x_{2} + 3 x_{4} - 2 x_{5} &= 1\\
- x_{2} - 3 x_{3} + 3 x_{4} - 2 x_{5} &= -11\\
- 3 x_{1} - 2 x_{2} + 2 x_{3} - x_{4} + 3 x_{5} &= 17\\
- 3 x_{2} + 3 x_{3} - 3 x_{5} &= 0
\end{aligned}
\]
\end{frame}
\begin{frame}[t]{Banco Gauss--Jordan: ejercicio 24 (5$\times$5)}
\textbf{Nivel:} Avanzado.\\
\textbf{Objetivo:} Ejecución exhaustiva del algoritmo.
\[
\begin{aligned}
3 x_{1} + 3 x_{3} + 3 x_{4} - 3 x_{5} &= 30\\
x_{1} + 3 x_{2} + x_{3} + 2 x_{4} + 3 x_{5} &= 18\\
2 x_{1} + 2 x_{2} + 3 x_{3} + 3 x_{5} &= 18\\
- 3 x_{1} + 3 x_{2} - 3 x_{3} + 2 x_{4} + 3 x_{5} &= -6\\
- 3 x_{1} - 2 x_{2} - 3 x_{3} + 2 x_{4} + 2 x_{5} &= -20
\end{aligned}
\]
\end{frame}
\begin{frame}[t]{Banco Gauss--Jordan: ejercicio 25 (5$\times$5)}
\textbf{Nivel:} Avanzado.\\
\textbf{Objetivo:} Ejecución exhaustiva del algoritmo.
\[
\begin{aligned}
x_{1} - 3 x_{2} - 2 x_{3} + 2 x_{4} - x_{5} &= -5\\
3 x_{1} + x_{2} + 3 x_{3} + 3 x_{4} - 2 x_{5} &= -14\\
- x_{1} - 2 x_{2} - 2 x_{3} + x_{4} + x_{5} &= 7\\
- 2 x_{1} + x_{2} + 3 x_{3} + 2 x_{4} + x_{5} &= 9\\
3 x_{1} - 2 x_{2} + 2 x_{3} - 2 x_{4} - 3 x_{5} &= -24
\end{aligned}
\]
\end{frame}
\begin{frame}[t]{Banco Gauss--Jordan: ejercicio 26 (5$\times$5)}
\textbf{Nivel:} Avanzado.\\
\textbf{Objetivo:} Ejecución exhaustiva del algoritmo.
\[
\begin{aligned}
- 2 x_{3} - x_{4} - 3 x_{5} &= 11\\
3 x_{2} + 3 x_{3} + x_{4} - x_{5} &= 4\\
x_{1} - 3 x_{2} - x_{3} - 3 x_{4} + x_{5} &= -17\\
2 x_{5} &= -6\\
3 x_{1} - 2 x_{3} - 3 x_{4} - 2 x_{5} &= 3
\end{aligned}
\]
\end{frame}
\begin{frame}[t]{Banco Gauss--Jordan: ejercicio 27 (5$\times$5)}
\textbf{Nivel:} Avanzado.\\
\textbf{Objetivo:} Ejecución exhaustiva del algoritmo.
\[
\begin{aligned}
- x_{1} - x_{2} + x_{3} + x_{4} + x_{5} &= 6\\
- x_{2} + x_{3} - x_{4} &= -4\\
2 x_{2} - 2 x_{3} + 2 x_{4} - x_{5} &= 4\\
3 x_{1} + 2 x_{2} + x_{3} - 2 x_{5} &= -12\\
x_{1} - 2 x_{2} + x_{3} - x_{4} - x_{5} &= -7
\end{aligned}
\]
\end{frame}
\begin{frame}[t]{Banco Gauss--Jordan: ejercicio 28 (5$\times$5)}
\textbf{Nivel:} Avanzado.\\
\textbf{Objetivo:} Ejecución exhaustiva del algoritmo.
\[
\begin{aligned}
- x_{1} - x_{2} - x_{3} + x_{4} - 2 x_{5} &= -6\\
- 2 x_{2} + 2 x_{3} - 3 x_{4} &= 2\\
3 x_{1} - x_{2} + x_{4} - 2 x_{5} &= 1\\
- 2 x_{1} + 2 x_{2} - 3 x_{3} - 2 x_{5} &= -1\\
- 2 x_{1} - x_{2} + 3 x_{3} + 2 x_{4} - 2 x_{5} &= -14
\end{aligned}
\]
\end{frame}
\begin{frame}[t]{Banco Gauss--Jordan: ejercicio 29 (3$\times$3)}
\textbf{Nivel:} Intermedio/Avanzado.\\
\textbf{Objetivo:} Variables libres e infinitas soluciones.
\[
\begin{aligned}
- x_{2} - 3 x_{3} &= 4\\
- 3 x_{1} + x_{2} + 3 x_{3} &= -16\\
- 3 x_{1} + x_{2} + 3 x_{3} &= -16
\end{aligned}
\]
\end{frame}
\begin{frame}[t]{Banco Gauss--Jordan: ejercicio 30 (3$\times$3)}
\textbf{Nivel:} Intermedio/Avanzado.\\
\textbf{Objetivo:} Variables libres e infinitas soluciones.
\[
\begin{aligned}
2 x_{1} + 3 x_{2} &= 20\\
x_{1} - 2 x_{2} - 3 x_{3} &= -4\\
x_{1} - 2 x_{2} - 3 x_{3} &= -4
\end{aligned}
\]
\end{frame}
\begin{frame}[t]{Banco Gauss--Jordan: ejercicio 31 (4$\times$4)}
\textbf{Nivel:} Intermedio/Avanzado.\\
\textbf{Objetivo:} Variables libres e infinitas soluciones.
\[
\begin{aligned}
- 3 x_{2} + 2 x_{3} - 2 x_{4} &= -17\\
3 x_{1} + 3 x_{2} - x_{3} - 3 x_{4} &= 0\\
2 x_{1} - 2 x_{4} &= -4\\
2 x_{1} - 2 x_{4} &= -4
\end{aligned}
\]
\end{frame}
\begin{frame}[t]{Banco Gauss--Jordan: ejercicio 32 (4$\times$4)}
\textbf{Nivel:} Intermedio/Avanzado.\\
\textbf{Objetivo:} Variables libres e infinitas soluciones.
\[
\begin{aligned}
x_{1} + 3 x_{2} - 3 x_{3} + 2 x_{4} &= -13\\
- 2 x_{2} + 2 x_{4} &= -4\\
- 2 x_{1} + 2 x_{2} - x_{3} + 3 x_{4} &= -17\\
- 2 x_{1} + 2 x_{2} - x_{3} + 3 x_{4} &= -17
\end{aligned}
\]
\end{frame}
\begin{frame}[t]{Banco Gauss--Jordan: ejercicio 33 (5$\times$5)}
\textbf{Nivel:} Intermedio/Avanzado.\\
\textbf{Objetivo:} Variables libres e infinitas soluciones.
\[
\begin{aligned}
2 x_{1} + 2 x_{2} - 3 x_{3} + x_{4} + x_{5} &= 9\\
- 3 x_{1} - 3 x_{2} + x_{3} + 3 x_{4} + 2 x_{5} &= 8\\
- x_{1} - x_{2} - x_{4} - 3 x_{5} &= 3\\
- x_{1} - x_{2} - 2 x_{3} - 3 x_{4} &= -5\\
- x_{1} - x_{2} - 2 x_{3} - 3 x_{4} &= -5
\end{aligned}
\]
\end{frame}
\begin{frame}[t]{Banco Gauss--Jordan: ejercicio 34 (5$\times$5)}
\textbf{Nivel:} Intermedio/Avanzado.\\
\textbf{Objetivo:} Variables libres e infinitas soluciones.
\[
\begin{aligned}
- 3 x_{1} - 2 x_{2} + 2 x_{3} - x_{4} + 2 x_{5} &= 7\\
2 x_{1} + x_{2} - 3 x_{3} - 2 x_{5} &= -3\\
- x_{1} + 2 x_{2} + 2 x_{4} - 2 x_{5} &= -10\\
2 x_{1} + x_{2} &= -2\\
2 x_{1} + x_{2} &= -2
\end{aligned}
\]
\end{frame}
\begin{frame}[t]{Banco Gauss--Jordan: ejercicio 35 (3$\times$3)}
\textbf{Nivel:} Intermedio/Avanzado.\\
\textbf{Objetivo:} Detección de incompatibilidad.
\[
\begin{aligned}
- 2 x_{1} + 2 x_{2} &= 8\\
3 x_{1} + x_{2} + 3 x_{3} &= 8\\
3 x_{1} + x_{2} + 3 x_{3} &= 9
\end{aligned}
\]
\end{frame}
\begin{frame}[t]{Banco Gauss--Jordan: ejercicio 36 (3$\times$3)}
\textbf{Nivel:} Intermedio/Avanzado.\\
\textbf{Objetivo:} Detección de incompatibilidad.
\[
\begin{aligned}
x_{2} - 3 x_{3} &= 2\\
3 x_{1} - 2 x_{2} - 3 x_{3} &= -10\\
3 x_{1} - 2 x_{2} - 3 x_{3} &= -9
\end{aligned}
\]
\end{frame}
\begin{frame}[t]{Banco Gauss--Jordan: ejercicio 37 (4$\times$4)}
\textbf{Nivel:} Intermedio/Avanzado.\\
\textbf{Objetivo:} Detección de incompatibilidad.
\[
\begin{aligned}
2 x_{2} + 3 x_{3} + 2 x_{4} &= -1\\
- 2 x_{1} + 3 x_{2} - 2 x_{3} &= -12\\
- 3 x_{2} - x_{3} - 2 x_{4} &= 5\\
- 3 x_{2} - x_{3} - 2 x_{4} &= 6
\end{aligned}
\]
\end{frame}
\begin{frame}[t]{Banco Gauss--Jordan: ejercicio 38 (4$\times$4)}
\textbf{Nivel:} Intermedio/Avanzado.\\
\textbf{Objetivo:} Detección de incompatibilidad.
\[
\begin{aligned}
3 x_{1} + x_{2} - x_{3} + x_{4} &= -2\\
x_{1} + 3 x_{2} + 2 x_{3} &= 4\\
- 2 x_{1} - 2 x_{2} + 2 x_{3} &= -10\\
- 2 x_{1} - 2 x_{2} + 2 x_{3} &= -9
\end{aligned}
\]
\end{frame}
\begin{frame}[t]{Banco Gauss--Jordan: ejercicio 39 (5$\times$5)}
\textbf{Nivel:} Intermedio/Avanzado.\\
\textbf{Objetivo:} Detección de incompatibilidad.
\[
\begin{aligned}
- 2 x_{1} + 3 x_{2} - 2 x_{3} + 2 x_{4} + x_{5} &= 6\\
- x_{1} - x_{2} + x_{3} + 2 x_{4} - 2 x_{5} &= 5\\
- x_{1} - 2 x_{2} + 3 x_{3} &= 10\\
- 3 x_{1} + 2 x_{2} - 2 x_{3} + x_{4} - 3 x_{5} &= -11\\
- 3 x_{1} + 2 x_{2} - 2 x_{3} + x_{4} - 3 x_{5} &= -10
\end{aligned}
\]
\end{frame}
\begin{frame}[t]{Banco Gauss--Jordan: ejercicio 40 (5$\times$5)}
\textbf{Nivel:} Intermedio/Avanzado.\\
\textbf{Objetivo:} Detección de incompatibilidad.
\[
\begin{aligned}
x_{1} - x_{2} - 3 x_{4} &= -7\\
- 2 x_{1} + 2 x_{2} + 3 x_{3} - 3 x_{4} &= 5\\
- x_{1} - 3 x_{3} - 2 x_{4} - 2 x_{5} &= -20\\
- 2 x_{1} - 3 x_{2} - 3 x_{3} + x_{4} - 2 x_{5} &= -21\\
- 2 x_{1} - 3 x_{2} - 3 x_{3} + x_{4} - 2 x_{5} &= -20
\end{aligned}
\]
\end{frame}
\begin{frame}[t]{Clave rápida: ejercicios 1--10}\begin{itemize}
\item E1: única, \(x_{1}=3, x_{2}=1\).
\item E2: única, \(x_{1}=2, x_{2}=4\).
\item E3: única, \(x_{1}=3, x_{2}=4\).
\item E4: única, \(x_{1}=-3, x_{2}=-1\).
\item E5: única, \(x_{1}=1, x_{2}=2\).
\item E6: única, \(x_{1}=4, x_{2}=4, x_{3}=2\).
\item E7: única, \(x_{1}=1, x_{2}=0, x_{3}=2\).
\item E8: única, \(x_{1}=-2, x_{2}=-2, x_{3}=4\).
\item E9: única, \(x_{1}=1, x_{2}=3, x_{3}=1\).
\item E10: única, \(x_{1}=-2, x_{2}=-3, x_{3}=-1\).
\end{itemize}\end{frame}
\begin{frame}[t]{Clave rápida: ejercicios 11--20}\begin{itemize}
\item E11: única, \(x_{1}=1, x_{2}=-1, x_{3}=2\).
\item E12: única, \(x_{1}=2, x_{2}=-2, x_{3}=3\).
\item E13: única, \(x_{1}=0, x_{2}=1, x_{3}=-1\).
\item E14: única, \(x_{1}=-1, x_{2}=1, x_{3}=0, x_{4}=-1\).
\item E15: única, \(x_{1}=-2, x_{2}=0, x_{3}=-3, x_{4}=4\).
\item E16: única, \(x_{1}=1, x_{2}=0, x_{3}=2, x_{4}=2\).
\item E17: única, \(x_{1}=-3, x_{2}=3, x_{3}=-3, x_{4}=0\).
\item E18: única, \(x_{1}=0, x_{2}=4, x_{3}=2, x_{4}=-1\).
\item E19: única, \(x_{1}=3, x_{2}=-2, x_{3}=3, x_{4}=2\).
\item E20: única, \(x_{1}=2, x_{2}=-2, x_{3}=3, x_{4}=-3\).
\end{itemize}\end{frame}
\begin{frame}[t]{Clave rápida: ejercicios 21--30}\begin{itemize}
\item E21: única, \(x_{1}=-3, x_{2}=3, x_{3}=3, x_{4}=-1, x_{5}=-3\).
\item E22: única, \(x_{1}=2, x_{2}=-1, x_{3}=2, x_{4}=-2, x_{5}=-1\).
\item E23: única, \(x_{1}=0, x_{2}=0, x_{3}=4, x_{4}=3, x_{5}=4\).
\item E24: única, \(x_{1}=3, x_{2}=3, x_{3}=3, x_{4}=3, x_{5}=-1\).
\item E25: única, \(x_{1}=-3, x_{2}=1, x_{3}=-1, x_{4}=1, x_{5}=3\).
\item E26: única, \(x_{1}=1, x_{2}=2, x_{3}=-3, x_{4}=4, x_{5}=-3\).
\item E27: única, \(x_{1}=0, x_{2}=-1, x_{3}=-2, x_{4}=3, x_{5}=4\).
\item E28: única, \(x_{1}=2, x_{2}=1, x_{3}=-1, x_{4}=-2, x_{5}=1\).
\item E29: infinitas, rango \(=2<3\).
\item E30: infinitas, rango \(=2<3\).
\end{itemize}\end{frame}
\begin{frame}[t]{Clave rápida: ejercicios 31--40}\begin{itemize}
\item E31: infinitas, rango \(=3<4\).
\item E32: infinitas, rango \(=3<4\).
\item E33: infinitas, rango \(=4<5\).
\item E34: infinitas, rango \(=4<5\).
\item E35: incompatible.
\item E36: incompatible.
\item E37: incompatible.
\item E38: incompatible.
\item E39: incompatible.
\item E40: incompatible.
\end{itemize}\end{frame}
\section{Ejercicios propuestos}
\begin{frame}[allowframebreaks,t]{Ejercicios propuestos}
\par\medskip{\large\bfseries A. Conceptos}\par\smallskip
\begin{enumerate}
\item Determine cuáles de las siguientes ecuaciones son lineales:
\[
3x-2y+z=4,\qquad x^2+y=1,\qquad xy-z=0,\qquad 2x_1-5x_2=7.
\]
\item Explique con sus palabras la diferencia entre un sistema compatible determinado, compatible
indeterminado e incompatible.
\item ¿Por qué todo sistema homogéneo es compatible?
\item ¿Qué información proporciona un pivote?
\item Explique la diferencia entre forma escalonada y forma escalonada reducida.
\end{enumerate}

\par\medskip{\large\bfseries B. Eliminación de Gauss}\par\smallskip
Resuelva y clasifique:
\begin{enumerate}
\item
\[
\begin{cases}
x+y=7,\\
2x-y=2.
\end{cases}
\]
\item
\[
\begin{cases}
x+y+z=6,\\
2x-y+3z=9,\\
3x+2y-z=4.
\end{cases}
\]
\item
\[
\begin{cases}
x+2y-z=3,\\
2x+4y-2z=6.
\end{cases}
\]
\item
\[
\begin{cases}
x-y=1,\\
2x-2y=5.
\end{cases}
\]
\end{enumerate}

\par\medskip{\large\bfseries C. Gauss--Jordan y variables libres}\par\smallskip
\begin{enumerate}
\item Reduzca a forma escalonada reducida:
\[
\left[
\begin{array}{ccc|c}
1&2&-1&3\\
2&5&1&8\\
-1&-1&2&-1
\end{array}
\right].
\]
\item Para
\[
\left[
\begin{array}{cccc|c}
1&2&0&-1&4\\
0&0&1&3&2\\
0&0&0&0&0
\end{array}
\right],
\]
identifique variables básicas y libres y describa todas las soluciones.
\end{enumerate}

\par\medskip{\large\bfseries D. Parámetros}\par\smallskip
Clasifique los sistemas según el valor del parámetro:
\begin{enumerate}
\item
\[
\begin{cases}
x+y=1,\\
2x+ay=2.
\end{cases}
\]
\item
\[
\begin{cases}
x-y=3,\\
2x-2y=k.
\end{cases}
\]
\end{enumerate}

\par\medskip{\large\bfseries E. Aplicaciones}\par\smallskip
\begin{enumerate}
\item Una mezcla de $20$ kg debe contener $35\%$ de un componente. ¿Cuántos kilogramos de mezclas
al $20\%$ y al $50\%$ deben combinarse?
\item Formule y resuelva un sistema para determinar dos números cuya suma es $18$ y cuya diferencia
es $4$.
\item Tres corrientes satisfacen
\[
I_1-I_2-I_3=0,\qquad 5I_1+2I_2=15,\qquad 2I_2-4I_3=0.
\]
Determine $I_1,I_2,I_3$.
\end{enumerate}
\end{frame}

\section{Resumen de la unidad}
\begin{frame}[allowframebreaks,t]{Resumen de la unidad}
Las ideas fundamentales pueden organizarse así:
\[
A\mathbf x=\mathbf b
\]
representa un sistema lineal. Las operaciones elementales producen sistemas equivalentes y permiten
reducir la matriz aumentada.

Durante la reducción:
\[
[0\ \cdots\ 0\mid c],\quad c\neq0
\]
indica incompatibilidad. Si no aparece contradicción, el sistema es compatible. Si todas las variables
son básicas, la solución es única; si existen variables libres, hay infinitas soluciones.

En términos de rango:
\[
\boxed{\rank(A)=\rank([A\mid\mathbf b])}
\]
es la condición de consistencia. Para un sistema compatible con $n$ incógnitas:
\[
\rank(A)=n\Longrightarrow\text{solución única},
\qquad
\rank(A)<n\Longrightarrow\text{infinitas soluciones}.
\]

Los métodos de matriz inversa y Cramer son herramientas adicionales bajo hipótesis específicas, mientras
que la eliminación gaussiana constituye el método general que articula gran parte de la unidad.
\end{frame}
\end{document}
